% !TEX program = pdflatex
\documentclass[11pt,a4paper]{article}

% =========================================================================
% EXPEDIENTE DEL PROYECTO
% Proyecto: CasaScore (AlquilaPuno)
% Gestión de Proyectos TICS (EST315) -- Semestre 2026-II -- Team Sobrevivientes
% =========================================================================

% --- PAQUETES DE CONFIGURACIÓN ---
\usepackage[utf8]{inputenc}
\usepackage[T1]{fontenc}
\usepackage[spanish,es-tabla]{babel}
\usepackage{geometry}
\geometry{
    top=3.6cm,
    bottom=2.3cm,
    left=2.2cm,
    right=2.2cm,
    % PARTE MODIFICABLE: headheight debe ser mayor que la altura del logo más alto (2.56cm = 73pt).
    headheight=80pt,
    headsep=15pt,
    footskip=20pt
}

\usepackage{graphicx}
\usepackage{fancyhdr}
\usepackage[table]{xcolor}
\usepackage{titlesec}
\usepackage{booktabs}
\usepackage{tabularx}
\usepackage{array}
\usepackage{amsmath}
\usepackage{enumitem}
\usepackage{microtype}
\usepackage{amssymb}
% --- RECURSOS VISUALES: GRÁFICOS Y CÓDIGO DE COLORES ---
\usepackage{tikz}
\usetikzlibrary{positioning,calc}
\definecolor{GrisFila}{RGB}{246, 248, 251}
\definecolor{CuadCerca}{RGB}{150, 30, 30}
\definecolor{CuadSatisf}{RGB}{180, 130, 20}
\definecolor{CuadInfo}{RGB}{10, 45, 90}
\definecolor{CuadMonit}{RGB}{110, 115, 125}
\newcommand{\etiqueta}[2]{{\setlength{\fboxsep}{2pt}\colorbox{#1}{\textcolor{white}{\scriptsize\bfseries #2}}}}
\newcommand{\cerca}{\etiqueta{CuadCerca}{Gestionar de cerca}}
\newcommand{\satisf}{\etiqueta{CuadSatisf}{Mantener satisfecho}}
\newcommand{\informado}{\etiqueta{CuadInfo}{Mantener informado}}
\newcommand{\monit}{\etiqueta{CuadMonit}{Monitorear}}
\usepackage{needspace}
\usepackage{longtable}
\definecolor{Rama1}{RGB}{10, 45, 90}
\definecolor{Rama2}{RGB}{40, 95, 150}
\definecolor{Rama3}{RGB}{180, 130, 20}
\definecolor{Rama4}{RGB}{30, 110, 60}
\definecolor{Rama5}{RGB}{110, 60, 140}
\definecolor{RaciR}{RGB}{150, 30, 30}
\definecolor{RaciA}{RGB}{10, 45, 90}
\definecolor{RaciC}{RGB}{180, 130, 20}
\definecolor{RaciI}{RGB}{140, 145, 155}
\definecolor{NivAlto}{RGB}{30, 110, 60}
\definecolor{NivMedio}{RGB}{180, 130, 20}
\definecolor{NivBajo}{RGB}{150, 30, 30}
\usetikzlibrary{patterns}
\definecolor{SevBaja}{RGB}{30, 110, 60}
\definecolor{SevMedia}{RGB}{180, 130, 20}
\definecolor{SevAlta}{RGB}{150, 30, 30}
\usetikzlibrary{arrows.meta,shapes.geometric,fit,backgrounds}
\usepackage{listings}
\lstset{basicstyle=\ttfamily\scriptsize, frame=single, rulecolor=\color{AzulUNA!40}, backgroundcolor=\color{GrisFila}, columns=fullflexible, keepspaces=true, breaklines=true, xleftmargin=4pt, xrightmargin=4pt, literate={á}{{\'a}}1 {é}{{\'e}}1 {í}{{\'i}}1 {ó}{{\'o}}1 {ú}{{\'u}}1 {ñ}{{\~n}}1 {Á}{{\'A}}1}

\usepackage{pdflscape}
\usepackage{adjustbox}
\usepackage{afterpage}
\usetikzlibrary{shapes.multipart,shapes.misc,decorations.pathreplacing}
\usepackage{hyperref}
\AtBeginDocument{\shorthandoff{<>}}
% --- HOJAS HORIZONTALES: encabezado y pie girados para leerse con la hoja apaisada ---
\fancypagestyle{horizontal}{%
  \fancyhf{}\renewcommand{\headrulewidth}{0pt}\renewcommand{\footrulewidth}{0pt}%
  \fancyhead[C]{\begin{tikzpicture}[remember picture,overlay]
    \node[rotate=90, anchor=center] at ($(current page.west)+(1.25cm,0.75cm)$) {%
      \begin{tabular}{@{}c@{}}\begin{minipage}{24.6cm}
      \raisebox{-0.5\height}{\includegraphics[height=1.15cm,keepaspectratio]{logo_fiei.png}}\hfill
      \begin{tabular}{c}{\scriptsize\bfseries\textcolor{AzulUNA}{UNIVERSIDAD NACIONAL DEL ALTIPLANO DE PUNO}}\\[-1pt]
      {\scriptsize\bfseries\textcolor{AzulUNA}{Gestión de Proyectos TICS (EST315) \textbar\ Team Sobrevivientes}}\\[-1pt]
      {\footnotesize\bfseries\textcolor{DoradoFIEI}{EXPEDIENTE DEL PROYECTO}}\end{tabular}\hfill
      \raisebox{-0.5\height}{\includegraphics[height=2.0cm,keepaspectratio]{logo_equipo.png}}
      \end{minipage}\end{tabular}};
    \node[rotate=90, anchor=center] at ($(current page.east)+(-1.0cm,0.75cm)$) {%
      \begin{minipage}{24.6cm}{\scriptsize\textcolor{GrisOscuro}{CasaScore (AlquilaPuno)}}\hfill{\small\thepage}\hfill{\scriptsize\textcolor{GrisOscuro}{Puno, Perú -- Semestre 2026-II}}\end{minipage}};
  \end{tikzpicture}}}
\newenvironment{hojah}[1]{\clearpage\begin{landscape}\thispagestyle{horizontal}\section{#1}}{\end{landscape}}
\newenvironment{grande}[1][13.6cm]{\par\noindent\begin{adjustbox}{width=\linewidth, max totalheight=#1, center}}{\end{adjustbox}\par}
\newcommand{\tile}[6]{\node[draw=#3, fill=#3!8, line width=1pt, rounded corners=4pt, minimum width=3.95cm, minimum height=2.0cm] at (#1,#2) {};
\node[font=\scriptsize\bfseries, text=GrisOscuro] at (#1,#2+0.6) {#4};
\node[font=\Large\bfseries, text=#3] at (#1,#2) {#5};
\node[font=\tiny, text=GrisOscuro] at (#1,#2-0.6) {#6};}
\tikzset{
 fl/.style={>=Stealth, font=\normalsize, node distance=0.9cm and 1.0cm},
 term/.style={rounded rectangle, draw=AzulUNA, fill=AzulUNA, text=white, font=\normalsize\bfseries, minimum height=0.9cm, inner xsep=10pt},
 proc/.style={rectangle, rounded corners=3pt, draw=AzulUNA, fill=AzulClaro, line width=0.9pt, align=center, minimum height=1.0cm, text width=#1, inner sep=4pt},
 proc/.default=3.4cm,
 dec/.style={diamond, aspect=2.1, draw=DoradoFIEI, fill=DoradoFIEI!12, line width=0.9pt, align=center, inner sep=1pt, font=\small\bfseries, text width=#1},
 dec/.default=2.4cm,
 dato/.style={trapezium, trapezium left angle=72, trapezium right angle=108, draw=Rama4, fill=Rama4!10, line width=0.9pt, align=center, text width=#1, minimum height=1.0cm},
 dato/.default=3.2cm,
 fle/.style={->, line width=1.1pt, AzulUNA},
 etq/.style={font=\small\bfseries, text=GrisOscuro, fill=white, inner sep=1.5pt},
 actor/.style={draw=#1, fill=#1!10, line width=1.2pt, rounded corners=6pt, align=center, font=\large},
}
\newcommand{\caja}[3]{\node[draw=#1, fill=#1!10, line width=1pt, rounded corners=4pt, align=center, inner sep=5pt] #2 {#3};}


% Logos: funciona a 2 o 3 niveles de profundidad dentro de docs/latex/
\graphicspath{{../../assets/img/}{../../../assets/img/}}

% --- PALETA INSTITUCIONAL FIEI / UNA PUNO ---
\definecolor{AzulUNA}{RGB}{10, 45, 90}
\definecolor{DoradoFIEI}{RGB}{180, 130, 20}
\definecolor{GrisOscuro}{RGB}{50, 50, 50}
\definecolor{AzulClaro}{RGB}{228, 236, 246}

% --- CONFIGURACIÓN DE LISTAS ---
\setlist{topsep=2pt, itemsep=2pt, parsep=1pt}

% --- COLUMNAS DE TABLA ---
\newcolumntype{L}{>{\raggedright\arraybackslash}X}
\newcolumntype{P}[1]{>{\raggedright\arraybackslash}p{#1}}
\newcolumntype{C}[1]{>{\centering\arraybackslash}p{#1}}
\newcolumntype{Q}[1]{>{\raggedleft\arraybackslash}p{#1}}
\newcolumntype{R}{>{\raggedleft\arraybackslash}X}
\renewcommand{\arraystretch}{1.25}

% --- FORMATO DE ENLACES ---
\hypersetup{
    colorlinks=true,
    linkcolor=AzulUNA,
    urlcolor=AzulUNA,
    citecolor=AzulUNA,
    pdftitle={Expediente del Proyecto -- CasaScore (AlquilaPuno)},
    pdfauthor={Team Sobrevivientes}
}

% --- FORMATO DE TÍTULOS ---
\titleformat{\section}{\color{AzulUNA}\normalfont\large\bfseries}{\thesection.}{0.7em}{}[\color{AzulUNA}\titlerule]
\titleformat{\subsection}{\color{AzulUNA}\normalfont\normalsize\bfseries}{\thesubsection.}{0.5em}{}
\titlespacing*{\section}{0pt}{9pt}{4pt}
\titlespacing*{\subsection}{0pt}{6pt}{2pt}

% --- ENCABEZADO Y PIE DE PÁGINA (fancyhdr) ---
\pagestyle{fancy}
\fancyhf{}

% ===== TAMAÑO DE LOS LOGOS (PARTE MODIFICABLE) =====
% logo_fiei.png  : height=1.45cm (lado izquierdo)
% logo_equipo.png: height=2.56cm (lado derecho; 25 % más grande porque la imagen tiene más borde blanco)
% Para cambiar el tamaño, modifique solo el valor height=... de cada \includegraphics.
% Si el logo más alto supera 2.6cm, aumente también headheight y top en la configuración de página.
\fancyhead[L]{%
    \includegraphics[height=1.45cm,keepaspectratio]{logo_fiei.png}%
}
\fancyhead[C]{%
    \scriptsize \textbf{\textcolor{AzulUNA}{UNIVERSIDAD NACIONAL DEL ALTIPLANO DE PUNO}}\\[2pt]
    \tiny \textcolor{GrisOscuro}{FACULTAD DE INGENIERÍA ESTADÍSTICA E INFORMÁTICA}\\[2pt]
    \scriptsize \textbf{\textcolor{AzulUNA}{Gestión de Proyectos TICS (EST315) \textbar\ Team Sobrevivientes}}\\[2pt]
    \footnotesize \textbf{\textcolor{DoradoFIEI}{EXPEDIENTE DEL PROYECTO}}%
}
\fancyhead[R]{%
    \includegraphics[height=2.56cm,keepaspectratio]{logo_equipo.png}%
}

\renewcommand{\headrulewidth}{0pt}
\renewcommand{\headrule}{}

\fancyfoot[L]{\scriptsize \textcolor{GrisOscuro}{CasaScore (AlquilaPuno)}}
\fancyfoot[C]{\small \thepage}
\fancyfoot[R]{\scriptsize \textcolor{GrisOscuro}{Puno, Perú -- Semestre 2026-II}}
\renewcommand{\footrulewidth}{0.4pt}

\begin{document}
\begin{center}
    {\Large \textbf{\textcolor{AzulUNA}{EXPEDIENTE DEL PROYECTO}}}\\[3pt]
    {\normalsize \textbf{CasaScore (AlquilaPuno)}}\\[4pt]
    \footnotesize \textbf{Fecha:} viernes 13/11/2026 (semana 12) \quad \textbar\quad \textbf{Versión:} 1.1 \quad \textbar\quad \textbf{Avance al corte:} viernes 06/11/2026
\end{center}

\section{El Proyecto en una Hoja}
{\small
\noindent\begin{tabularx}{\textwidth}{@{}P{3.4cm} L@{}}
\toprule
\textbf{Problema:} & En Puno, quien busca cuarto no sabe si el precio es justo, si el lugar es seguro ni qué reglas de convivencia tiene hasta que lo visita. \\
\textbf{Producto:} & App que ordena los cuartos en alquiler con un puntaje de 0 a 100, el \textbf{CASA SCORE}, según precio, cercanía al punto que elige el arrendatario, seguridad, servicios y tipo de baño, y muestra siempre el desglose. \\
\textbf{Clientes:} & Arrendatarios (estudiantes, trabajadores y familias) y propietarios de cuartos. \\
\textbf{Sponsor:} & Team Sobrevivientes, que financia el gasto, aprueba los gates y es dueño del producto. Asesor: Ing. Renzo Apaza Cutipa. \\
\textbf{Equipo:} & Clyde N. Paricahua Pari (Project Manager), Ruth K. Apaza Solis (Planificación), Eliseo Tarqui Ajahuana (Técnico) y Viviana M. Quisbert Quispe (Comunicación, Riesgos y Documentación). \\
\textbf{Duración y esfuerzo:} & 16 semanas (31/08 al 18/12/2026); 8 h por semana por integrante; 512 h en total. \\
\textbf{Enfoque:} & Híbrido: líneas base fijas de alcance, cronograma y costo; el software se construye en sprints de una semana con un tablero Scrum y Kanban. \\
\textbf{Repositorio y tablero:} & \texttt{github.com/xClaIDx/ProyectoTics} y \texttt{github.com/users/xClaIDx/projects/1} \\
\bottomrule
\end{tabularx}
}

\section{Cifras Clave}
\begin{center}
\begin{tikzpicture}
\tile{0}{0}{AzulUNA}{LÍNEA BASE (BAC)}{S/ 11,663.47}{presupuesto total S/ 12,829.82}
\tile{4.25}{0}{AzulUNA}{ESFUERZO}{512 h}{5 ramas y 23 paquetes}
\tile{8.5}{0}{AzulUNA}{REQUISITOS}{9 RF + 8 RNF}{RF-01 a RF-09 en el MVP}
\tile{12.75}{0}{AzulUNA}{PLAZO}{16 semanas}{75\,\% de terminar a tiempo}
\tile{0}{-2.3}{Rama4}{VAN ECONÓMICO}{S/ 65,781.59}{beneficio social a 3 años}
\tile{4.25}{-2.3}{Rama4}{PAYBACK}{0.92 años}{descontado: 1.01 años}
\tile{8.5}{-2.3}{Rama4}{TIR ECONÓMICA}{169.9\,\%}{relación B/C de 6.92}
\tile{12.75}{-2.3}{Rama4}{GASTO EN EFECTIVO}{S/ 2,507.44}{S/ 626.86 por integrante}
\end{tikzpicture}
\end{center}

\Needspace{16\baselineskip}
\section{Objetivos y Cómo se Miden}
{\small
\rowcolors{2}{white}{GrisFila}
\noindent\begin{tabularx}{\textwidth}{@{}C{0.7cm} L P{4.6cm} C{1.7cm} C{1.7cm}@{}}
\toprule
\rowcolor{AzulClaro}
\textbf{N.\textsuperscript{o}} & \textbf{Objetivo específico} & \textbf{Indicador y meta} & \textbf{Fecha} & \textbf{Estado} \\
\midrule
O1 & Entregar el MVP de la Unidad I con motor y ranking explicado. & Ranking en 250 ms o menos; 5 cuartos de prueba en el orden esperado; 7 de 7 pruebas del motor. & 09/10 & \etiqueta{Rama4}{Cumplido} \\
O2 & Entregar el MVP final con registro, punto de referencia y datos en la nube. & RF-01 a RF-09 aceptados; 0 errores críticos abiertos. & 04/12 & \etiqueta{SevMedia}{En curso} \\
O3 & Lograr la satisfacción de los arrendatarios del piloto. & 80\,\% o más de 10 arrendatarios califica 4 o 5 sobre 5. & 27/11 & \etiqueta{GrisOscuro}{Piloto} \\
O4 & Facilitar la publicación a los propietarios. & 5 propietarios publican en 10 min o menos; 80\,\% o más satisfecho. & 27/11 & \etiqueta{GrisOscuro}{Piloto} \\
O5 & Terminar dentro del tiempo y del presupuesto. & SPI y CPI de 0.95 o más; gasto en efectivo de S/ 2,507.44 o menos. & 18/12 & \etiqueta{SevMedia}{CPI 0.95} \\
\bottomrule
\end{tabularx}
}

\Needspace{30\baselineskip}
\section{Objetivos por Iteración}
Cada sprint dura una semana (sábado a viernes) y tiene un objetivo medible que se revisa el viernes de 20:00 a 21:00. Los sprints 1 a 7 ya se cerraron; desde el sprint 8 son el compromiso del equipo.

\vspace{0.1cm}
{\footnotesize
\setlength{\tabcolsep}{3.5pt}
\rowcolors{2}{white}{GrisFila}
\noindent\begin{tabularx}{\textwidth}{@{}C{0.9cm} C{0.8cm} L P{4.4cm} C{1.6cm}@{}}
\toprule
\rowcolor{AzulClaro}
\textbf{Sprint} & \textbf{Sem.} & \textbf{Objetivo del sprint} & \textbf{Se logra si} & \textbf{Resultado} \\
\midrule
1 & 4 & Terminar la planificación base y avanzar el motor CASA SCORE. & Motor al 45\,\%; planificación al 95\,\%. & \etiqueta{SevAlta}{SPI 0.78} \\
2 & 5 & Cerrar el motor y el modelo de datos; decidir sobre Horus Sentinel. & Motor con 4 criterios; acta de exclusión votada. & \etiqueta{SevAlta}{SPI 0.65} \\
3 & 6 & Terminar la publicación de cuartos e iniciar el levantamiento. & Formulario con 4 fotos; 5 cuartos reales. & \etiqueta{SevAlta}{SPI 0.56} \\
4 & 7 & Entregar el prototipo con ranking explicado. & Ranking de 5 cuartos en el orden esperado. & \etiqueta{Rama4}{Cumplido} \\
5 & 8 & Sustentar la Unidad I y cerrar ranking y detalle. & Sustentación aprobada; 9 cuartos. & \etiqueta{Rama4}{Cumplido} \\
6 & 10 & Aprobar el Plan de Dirección y preparar la nube. & Plan v1.1 firmado; proyecto Supabase creado. & \etiqueta{Rama4}{Cumplido} \\
7 & 11 & Avanzar el registro en la nube y medir el desempeño. & Registro al 60\,\%; corte de valor ganado. & \etiqueta{SevMedia}{40\,\%} \\
8 & 12 & Punto de referencia, convocatoria y cambio del piloto. & 10 inscritos; SC-02 votada; plan de cierre. & \etiqueta{AzulUNA}{En curso} \\
9 & 13 & Registro e inicio de sesión con Supabase y perfil. & Inicio de sesión con RLS (18/11); 12 cuartos. & \etiqueta{GrisOscuro}{Plan} \\
10 & 14 & Piloto con clientes reales. & 10 arrendatarios y 5 propietarios; encuesta. & \etiqueta{GrisOscuro}{Plan} \\
11 & 15 & Corregir defectos y aceptar el MVP final. & 0 errores críticos; 14 casos de prueba; v1.0. & \etiqueta{GrisOscuro}{Plan} \\
12 & 16 & Informe final y último corte de valor ganado. & Informe revisado; SPI y CPI finales. & \etiqueta{GrisOscuro}{Plan} \\
\bottomrule
\end{tabularx}
}

\vspace{0.1cm}
{\small\noindent La semana 9 (exámenes parciales) no tiene sprint y la semana 17 (exámenes finales) se dedica al cierre: exposición el 14/12 y acta de cierre el 18/12.}

\Needspace{14\baselineskip}
\section{Requisitos del MVP}
{\footnotesize
\rowcolors{2}{white}{GrisFila}
\noindent\begin{tabularx}{\textwidth}{@{}C{1.1cm} P{3.4cm} L C{1.8cm}@{}}
\toprule
\rowcolor{AzulClaro}
\textbf{Código} & \textbf{Requisito} & \textbf{Qué hace la app} & \textbf{Estado} \\
\midrule
RF-01 & Registro e inicio de sesión & Cuentas de arrendatario y propietario; cada rol ve sus opciones. & \etiqueta{SevMedia}{En curso} \\
RF-02 & Publicación de cuartos & Formulario en 5 pasos con 4 fotos obligatorias. & \etiqueta{Rama4}{Hecho} \\
RF-03 & Punto de referencia & Distancia y tiempo a pie desde el punto marcado en el mapa. & \etiqueta{Rama4}{Web} \\
RF-04 & Prioridades & Pesos de 0 a 5 para precio, cercanía, seguridad, servicios y baño. & \etiqueta{Rama4}{Hecho} \\
RF-05 & Normalización 0--100 & Cada criterio a una escala común (S/ 140 da 100; S/ 850 da 0). & \etiqueta{Rama4}{Hecho} \\
RF-06 & CASA SCORE y ranking & Promedio ponderado, bonificaciones y penalizaciones, de 0 a 100. & \etiqueta{Rama4}{Hecho} \\
RF-07 & Explicación del puntaje & Desglose de los 5 criterios, bonificaciones y penalizaciones. & \etiqueta{Rama4}{Hecho} \\
RF-08 & Perfil y mis cuartos & Prioridades guardadas; editar u ocultar cuartos. & \etiqueta{SevMedia}{En curso} \\
RF-09 & Datos en la nube & Usuarios, cuartos y preferencias compartidos en Supabase. & \etiqueta{SevMedia}{En curso} \\
\bottomrule
\end{tabularx}
}

\vspace{0.1cm}
{\footnotesize\noindent\textbf{No funcionales:} ranking en menos de 250 ms con 500 cuartos; Android 10 o superior; contraseñas cifradas, seguridad por fila y HTTPS; Ley N.\textsuperscript{o} 29733; publicar en 10 minutos o menos; 100\,\% de puntajes con desglose; servicios en plan gratuito; código revisado en cada Pull Request.}

\begin{hojah}{Diagrama Cliente--Servicio}
\begin{grande}
\begin{tikzpicture}[fl]
\node[actor=AzulUNA, minimum width=6.6cm, minimum height=3.2cm, font=\Large] (cs) at (0,0) {\textbf{CasaScore}\\[2pt]{\normalsize ranking de cuartos con}\\{\normalsize puntaje explicado de 0 a 100}};
\node[actor=Rama4, minimum width=4.6cm, minimum height=2.4cm] (arr) at (-11.6,0.0) {\textbf{Arrendatario}\\{\small estudiante, trabajador}\\{\small o familia}};
\node[actor=DoradoFIEI, minimum width=4.6cm, minimum height=2.4cm] (pro) at (11.6,0.0) {\textbf{Propietario}\\{\small dueño de uno o}\\{\small más cuartos}};
\node[actor=Rama5, minimum width=4.2cm, minimum height=1.5cm] (sb) at (-4.6,-5.0) {\textbf{Supabase}\\{\small cuentas, datos y fotos}};
\node[actor=GrisOscuro, minimum width=4.2cm, minimum height=1.5cm] (osm) at (0,-5.0) {\textbf{OpenStreetMap}\\{\small mapa de Puno}};
\node[actor=RaciR, minimum width=4.2cm, minimum height=1.5cm] (eq) at (4.6,-5.0) {\textbf{Team Sobrevivientes}\\{\small soporte y moderación}};
\draw[fle, Rama4] ([yshift=0.7cm]arr.east) -- node[above, font=\normalsize, align=center]{prioridades (pesos 0 a 5)\\y punto de referencia} ([yshift=0.7cm]cs.west);
\draw[fle, Rama4] ([yshift=-0.7cm]cs.west) -- node[below, font=\normalsize, align=center]{ranking, desglose del puntaje,\\4 fotos y contacto} ([yshift=-0.7cm]arr.east);
\draw[fle, DoradoFIEI!80!black] ([yshift=0.7cm]pro.west) -- node[above, font=\normalsize, align=center]{datos del cuarto, 4 fotos\\y condiciones de convivencia} ([yshift=0.7cm]cs.east);
\draw[fle, DoradoFIEI!80!black] ([yshift=-0.7cm]cs.east) -- node[below, font=\normalsize, align=center]{publicación visible e\\inquilinos interesados} ([yshift=-0.7cm]pro.west);
\draw[<->, line width=1.1pt, Rama5, >=Stealth] (cs.south west) -- node[etq, left=2pt]{HTTPS} (sb.north);
\draw[<->, line width=1.1pt, GrisOscuro, >=Stealth] (cs.south) -- node[etq, right=2pt]{teselas} (osm.north);
\draw[<->, line width=1.1pt, RaciR, >=Stealth] (cs.south east) -- node[etq, right=2pt]{mejoras} (eq.north);
\draw[<->, line width=1.1pt, dashed, GrisOscuro, >=Stealth] (arr.north) to[out=50, in=130] node[above, font=\normalsize, fill=white]{visita, acuerdo y pago del alquiler (fuera de la app)} (pro.north);
\end{tikzpicture}
\end{grande}
\vspace{0.2cm}
{\large\noindent\textbf{Propuesta de valor:} el arrendatario compara cuartos con un mismo criterio y conoce las reglas antes de visitar; el propietario publica en 10 minutos y llega a inquilinos que ya saben lo que ofrece. Ningún pago cambia la posición de un cuarto.}
\end{hojah}

\begin{hojah}{Mapa del Servicio: Buscar y Elegir un Cuarto}
\begin{grande}
\begin{tikzpicture}[font=\small]
\def\W{3.15}
\foreach \i/\t in {0/Ingresa,1/Elige prioridades,2/Marca su punto,3/Revisa el ranking,4/Abre el detalle,5/Contacta,6/Visita y alquila} {
  \node[draw=AzulUNA, fill=AzulUNA, text=white, font=\small\bfseries, text width=2.8cm, align=center, minimum height=0.9cm] at (2.6+\i*\W,0) {\t};}
\foreach \y/\l/\c in {-1.75/{Acciones del\\arrendatario}/Rama4, -4.15/{Lo que ve\\en la app}/AzulUNA, -6.55/{Procesos\\internos}/DoradoFIEI, -8.95/{Soporte}/Rama5} {
  \node[draw=\c, fill=\c!15, font=\normalsize\bfseries, text width=2.0cm, align=center, minimum height=2.1cm] at (0,\y) {\l};}
\newcommand{\cel}[3]{\node[draw=gray!50, fill=white, text width=2.85cm, align=center, minimum height=2.1cm] at (2.6+#1*\W,#2) {#3};}
\cel{0}{-1.75}{Crea su cuenta o inicia sesión}
\cel{1}{-1.75}{Mueve los pesos de 0 a 5}
\cel{2}{-1.75}{Marca la UNA, su trabajo u otro lugar}
\cel{3}{-1.75}{Compara puntajes y filtra}
\cel{4}{-1.75}{Mira 4 fotos y el desglose}
\cel{5}{-1.75}{Escribe al propietario}
\cel{6}{-1.75}{Acuerda y paga fuera de la app}
\cel{0}{-4.15}{Pantalla de ingreso por rol}
\cel{1}{-4.15}{Panel de prioridades y perfiles rápidos}
\cel{2}{-4.15}{Mapa con lugares de Puno}
\cel{3}{-4.15}{Lista ordenada con su CASA SCORE}
\cel{4}{-4.15}{Detalle con bonos, penalidades y mapa}
\cel{5}{-4.15}{Botón de WhatsApp}
\cel{6}{-4.15}{--}
\cel{0}{-6.55}{Valida correo y contraseña; crea la sesión}
\cel{1}{-6.55}{Guarda las preferencias}
\cel{2}{-6.55}{Haversine y tiempo a pie}
\cel{3}{-6.55}{Motor CASA SCORE ordena}
\cel{4}{-6.55}{Explica cada criterio}
\cel{5}{-6.55}{Registra el interés}
\cel{6}{-6.55}{--}
\cel{0}{-8.95}{Supabase Auth}
\cel{1}{-8.95}{Base de datos con seguridad por fila}
\cel{2}{-8.95}{OpenStreetMap}
\cel{3}{-8.95}{Cuartos activos del propietario}
\cel{4}{-8.95}{Fotos en el almacenamiento}
\cel{5}{-8.95}{Propietario}
\cel{6}{-8.95}{Propietario}
\draw[RaciR, line width=1.3pt, dashed] (-1.1,-2.95) -- (23.0,-2.95) node[right, font=\small\bfseries, text=RaciR, align=left]{línea de\\interacción};
\draw[RaciR, line width=1.3pt, dashed] (-1.1,-5.35) -- (23.0,-5.35) node[right, font=\small\bfseries, text=RaciR, align=left]{línea de\\visibilidad};
\foreach \i in {0,...,5} {\draw[->, >=Stealth, thick, AzulUNA] (2.6+\i*\W+1.5,0) -- (2.6+\i*\W+\W-1.5,0);}
\end{tikzpicture}
\end{grande}
\end{hojah}

\begin{hojah}{Diagrama de Flujo: el Arrendatario Busca un Cuarto}
\begin{grande}
\begin{tikzpicture}[fl]
\node[term] (ini) {Inicio};
\node[dec, right=of ini] (cta) {Tiene cuenta?};
\node[proc=3.0cm, above=1.0cm of cta] (reg) {Crear cuenta como arrendatario};
\node[proc=3.0cm, right=1.2cm of cta] (ses) {Iniciar sesión};
\node[dec, right=1.2cm of ses] (pri) {Tiene prioridades guardadas?};
\node[proc=3.2cm, above=1.0cm of pri] (pes) {Elegir pesos de 0 a 5 y marcar el punto};
\node[proc=3.2cm, right=1.2cm of pri] (cal) {Calcular el CASA SCORE de cada cuarto};
\node[proc=3.2cm, below=1.4cm of cal] (lis) {Ver la lista ordenada de mayor a menor};
\node[dec, left=1.2cm of lis] (int) {Le interesa un cuarto?};
\node[proc=3.2cm, left=1.2cm of int] (aju) {Ajustar filtros, pesos o punto};
\node[proc=3.2cm, below=1.4cm of int] (det) {Ver el detalle: 4 fotos y desglose};
\node[dec, right=1.2cm of det] (vis) {Quiere visitarlo?};
\node[proc=3.2cm, right=1.2cm of vis] (wsp) {Contactar al propietario por WhatsApp};
\node[term, below=1.0cm of wsp] (fin) {Fin};
\draw[fle] (ini) -- (cta);
\draw[fle] (cta) -- node[etq, above]{sí} (ses);
\draw[fle] (cta) -- node[etq, left]{no} (reg);
\draw[fle] (reg) -| (ses);
\draw[fle] (ses) -- (pri);
\draw[fle] (pri) -- node[etq, above]{sí} (cal);
\draw[fle] (pri) -- node[etq, left]{no} (pes);
\draw[fle] (pes) -| (cal);
\draw[fle] (cal) -- (lis);
\draw[fle] (lis) -- (int);
\draw[fle] (int) -- node[etq, above]{no} (aju);
\draw[fle] (aju) |- ([yshift=-0.35cm]cal.west);
\draw[fle] (int) -- node[etq, right]{sí} (det);
\draw[fle] (det) -- (vis);
\draw[fle] (vis) -- node[etq, above]{sí} (wsp);
\draw[fle] (vis.south) |- ++(0,-0.7) -| node[etq, pos=0.25, below]{no, volver a la lista} ([xshift=-0.6cm]lis.south);
\draw[fle] (wsp) -- (fin);
\end{tikzpicture}
\end{grande}
\end{hojah}

\begin{hojah}{Diagrama de Flujo: el Propietario Publica un Cuarto}
\begin{grande}
\begin{tikzpicture}[fl]
\node[term] (ini) {Inicio};
\node[proc=3.0cm, right=of ini] (ses) {Iniciar sesión como propietario};
\node[proc=3.0cm, right=of ses] (p1) {Paso 1: datos básicos\\{\small precio, área, baño}};
\node[proc=3.0cm, right=of p1] (p2) {Paso 2: servicios y seguridad};
\node[proc=3.0cm, right=of p2] (p3) {Paso 3: condiciones de convivencia};
\node[proc=3.0cm, below=1.6cm of p3] (p4) {Paso 4: ubicación en el mapa};
\node[proc=3.0cm, left=of p4] (p5) {Paso 5: 4 fotos\\{\small fachada, habitación, baño, ventana}};
\node[dec=2.6cm, left=of p5] (ok) {Datos obligatorios completos?};
\node[proc=3.0cm, below=1.4cm of ok] (fal) {Marcar los campos que faltan};
\node[proc=3.0cm, left=of ok] (pre) {Vista previa del puntaje};
\node[proc=3.0cm, below=1.4cm of pre] (pub) {Publicar: el cuarto entra al ranking};
\node[dec=2.4cm, right=4.6cm of pub, yshift=-2.6cm] (cam) {Cambió algo?};
\node[proc=3.0cm, right=1.2cm of cam] (mis) {Editar u ocultar en Mis cuartos};
\node[term, left=1.2cm of cam] (fin) {Fin};
\draw[fle] (ini) -- (ses);
\draw[fle] (ses) -- (p1);
\draw[fle] (p1) -- (p2);
\draw[fle] (p2) -- (p3);
\draw[fle] (p3) -- (p4);
\draw[fle] (p4) -- (p5);
\draw[fle] (p5) -- (ok);
\draw[fle] (ok) -- node[etq, above]{sí} (pre);
\draw[fle] (ok) -- node[etq, right]{no} (fal);
\draw[fle] (fal) -| (p5);
\draw[fle] (pre) -- (pub);
\draw[fle] (pub) |- (cam);
\draw[fle] (cam) -- node[etq, above]{sí} (mis);
\draw[fle] (cam) -- node[etq, above]{no} (fin);
\end{tikzpicture}
\end{grande}
\vspace{0.3cm}
{\large\noindent Sin un dato obligatorio no se puede avanzar al siguiente paso. El borrador se guarda solo, las fotos se comprimen antes de subirlas y ocultar las condiciones de convivencia resta 25 puntos.}
\end{hojah}

\begin{hojah}{Diagrama de Flujo: Cálculo del CASA SCORE}
\begin{grande}
\begin{tikzpicture}[fl]
\node[term] (ini) {Inicio};
\node[dato=3.6cm, right=of ini] (in) {Cuartos activos, 5 pesos $w_i$ y punto de referencia};
\node[proc=3.6cm, right=of in] (cu) {Tomar el siguiente cuarto};
\node[proc=3.8cm, right=of cu] (geo) {Distancia $d$ (Haversine) y tiempo $t = 1.3\,d/80$ min};
\node[proc=3.8cm, below=1.3cm of geo] (nor) {Normalizar 0--100: precio, cercanía, seguridad, servicios y baño};
\node[proc=3.6cm, left=of nor] (base) {Base $= \dfrac{\sum w_i N_i}{\sum w_i}$};
\node[proc=3.6cm, left=of base] (bon) {Sumar bonos: luz +8, entrada libre +5, 16 m$^2$ +4};
\node[proc=3.6cm, below=1.3cm of bon] (pen) {Restar penalidades; ocultar condiciones $-25$};
\node[proc=3.6cm, right=of pen] (lim) {Limitar a 0--100 y guardar el desglose};
\node[dec=2.4cm, right=of lim] (mas) {Quedan cuartos?};
\node[proc=3.4cm, right=of mas] (ord) {Ordenar de mayor a menor};
\node[term, below=0.9cm of ord] (fin) {Mostrar ranking};
\draw[fle] (ini) -- (in);
\draw[fle] (in) -- (cu);
\draw[fle] (cu) -- (geo);
\draw[fle] (geo) -- (nor);
\draw[fle] (nor) -- (base);
\draw[fle] (base) -- (bon);
\draw[fle] (bon) -- (pen);
\draw[fle] (pen) -- (lim);
\draw[fle] (lim) -- (mas);
\draw[fle] (mas) -- node[etq, above]{no} (ord);
\draw[fle] (mas.north) -- node[etq, right]{sí} ++(0,1.2) -| ([xshift=0.5cm]cu.south);
\draw[fle] (ord) -- (fin);
\end{tikzpicture}
\end{grande}
\vspace{0.3cm}
{\large\noindent Con 500 cuartos el cálculo toma pocos milisegundos (meta: 250 ms). El motor es el mismo en la app y en la API, y lo verifican 7 pruebas automáticas.}
\end{hojah}

\begin{hojah}{UML: Diagrama de Casos de Uso}
\begin{grande}
\begin{tikzpicture}[>=Stealth, cu/.style={ellipse, draw=AzulUNA, fill=AzulClaro, line width=1pt, align=center, font=\normalsize, minimum width=4.4cm, minimum height=1.2cm, inner sep=1pt}]
\newcommand{\monito}[3]{\begin{scope}[shift={(#1)}, line width=1.4pt, draw=#3]
\draw (0,0.9) circle (0.3); \draw (0,0.6) -- (0,-0.3); \draw (-0.5,0.3) -- (0.5,0.3); \draw (0,-0.3) -- (-0.4,-1.0); \draw (0,-0.3) -- (0.4,-1.0);
\node[font=\large\bfseries, text=#3, below] at (0,-1.05) {#2};\end{scope}}
\draw[line width=1.4pt, AzulUNA, rounded corners=8pt] (-6.8,-7.6) rectangle (6.8,6.4);
\node[font=\Large\bfseries, text=AzulUNA, anchor=north west] at (-6.6,6.3) {Sistema CasaScore};
\node[cu] (u1) at (0,5.0) {Registrarse e iniciar sesión};
\node[cu] (u2) at (-3.4,3.2) {Definir prioridades};
\node[cu] (u3) at (-3.4,1.6) {Marcar punto de referencia};
\node[cu] (u4) at (-3.4,0.0) {Ver ranking};
\node[cu, fill=DoradoFIEI!15, draw=DoradoFIEI] (u5) at (1.2,-1.3) {Calcular CASA SCORE};
\node[cu] (u6) at (-3.4,-2.2) {Ver detalle y desglose};
\node[cu] (u7) at (-3.4,-3.9) {Guardar favorito};
\node[cu] (u8) at (-3.4,-5.6) {Contactar por WhatsApp};
\node[cu] (u9) at (3.4,2.4) {Publicar cuarto};
\node[cu] (u10) at (3.4,0.6) {Subir 4 fotos};
\node[cu] (u11) at (3.4,-3.0) {Ver vista previa del puntaje};
\node[cu] (u12) at (3.4,-5.0) {Editar u ocultar cuarto};
\monito{-10.0,1.2}{Arrendatario}{Rama4}
\monito{10.0,-0.6}{Propietario}{DoradoFIEI!80!black}
\monito{10.0,4.6}{Supabase}{Rama5}
\foreach \u in {u1,u2,u3,u4,u6,u7,u8} {\draw[line width=0.9pt, Rama4] (-9.5,1.2) -- (\u.west);}
\foreach \u in {u1,u9,u12} {\draw[line width=0.9pt, DoradoFIEI!80!black] (9.5,-0.6) -- (\u.east);}
\draw[line width=0.9pt, Rama5] (9.5,4.6) -- (u1.east);
\draw[->, dashed, thick] (u4.east) -- node[etq, sloped, above, pos=0.55]{\guillemotleft include\guillemotright} (u5.north west);
\draw[->, dashed, thick] (u9) -- node[etq, right]{\guillemotleft include\guillemotright} (u10);
\draw[->, dashed, thick] (u11) -- node[etq, right]{\guillemotleft extend\guillemotright} (u9.south east);
\draw[->, dashed, thick] (u6.east) -- node[etq, sloped, below, pos=0.55]{\guillemotleft include\guillemotright} (u5.south west);
\end{tikzpicture}
\end{grande}
\end{hojah}

\begin{hojah}{UML: Diagrama de Clases}
\begin{grande}
\begin{tikzpicture}[>=Stealth, cl/.style={rectangle split, rectangle split parts=3, draw=#1, line width=1.1pt, rectangle split part fill={#1!25,white,white}, font=\normalsize, text width=5.0cm, align=left, inner sep=4pt}]
\node[cl=AzulUNA] (usu) at (0,4) {\centering\textbf{Usuario}\nodepart{second}id: uuid\\nombre: texto\\correo: texto\\rol: arrendatario \textbar\ propietario\nodepart{third}+ registrarse()\\+ iniciarSesion(): Sesion};
\node[cl=Rama4] (pref) at (-7.6,4) {\centering\textbf{Preferencia}\nodepart{second}pesos: precio, cercanía, seguridad, servicios, baño (0 a 5)\\refLat, refLng: decimal\nodepart{third}+ guardar()};
\node[cl=GrisOscuro] (ses) at (7.6,5.2) {\centering\textbf{Sesion}\nodepart{second}token: texto\\creadaEn: fecha\nodepart{third}+ cerrar()};
\node[cl=DoradoFIEI] (cua) at (0,-2.6) {\centering\textbf{Cuarto}\nodepart{second}id, titulo, direccion\\precio: S/, area: m$^2$\\banoPrivado, luzNatural: lógico\\lat, lng: decimal\\estado: activo \textbar\ oculto\nodepart{third}+ publicar()\\+ ocultar()};
\node[cl=Rama2] (serv) at (7.6,0.2) {\centering\textbf{Servicios}\nodepart{second}agua24h, luz24h, internet24h\nodepart{third}};
\node[cl=Rama2] (cond) at (7.6,-3.2) {\centering\textbf{Condiciones}\nodepart{second}entradaLibre, horaLlegada\\visitas, cocina, mascotas\\declaracionTransparente\nodepart{third}};
\node[cl=Rama5] (foto) at (7.6,-6.6) {\centering\textbf{Foto}\nodepart{second}tipo: fachada, habitación, baño, ventana\\url: texto\nodepart{third}};
\node[cl=RaciR] (mot) at (-7.6,-2.4) {\centering\textbf{MotorCasaScore}\nodepart{second}bonos y penalidades\nodepart{third}+ calcular(cuarto, pref): Desglose\\+ ordenar(cuartos): lista};
\node[cl=RaciR] (des) at (-7.6,-6.8) {\centering\textbf{Desglose}\nodepart{second}criterios: 5 valores 0 a 100\\bonos, penalidades\\total: 0 a 100\nodepart{third}};
\draw[thick] (usu.west) -- node[etq, above, pos=0.15]{1} node[etq, above, pos=0.85]{0..1} (pref.east);
\draw[thick] (usu.east) -- node[etq, above, pos=0.15]{1} node[etq, above, pos=0.85]{*} (ses.west);
\draw[thick] (usu.south) -- node[etq, right, pos=0.12]{1 propietario} node[etq, right, pos=0.82]{publica *} (cua.north);
\draw[thick, dashed] ([xshift=1.2cm]usu.south) to[out=-60,in=60] node[etq, right]{favorito * a *} ([xshift=1.2cm]cua.north);
\draw[thick] (cua.east) -- node[etq, above, pos=0.8]{1} (serv.west);
\draw[thick] (cua.east) -- node[etq, above, pos=0.8]{1} (cond.west);
\draw[thick] (cua.east) -- node[etq, below, pos=0.8]{4..*} (foto.west);
\draw[->, thick, dashed] (mot.east) -- node[etq, above]{usa} (cua.west);
\draw[->, thick, dashed] (mot.north) -- node[etq, right]{usa} (pref.south);
\draw[->, thick] (mot.south) -- node[etq, right]{produce} (des.north);
\end{tikzpicture}
\end{grande}
\end{hojah}

\begin{hojah}{UML: Diagrama de Secuencia del Ranking}
\begin{grande}
\begin{tikzpicture}[>=Stealth, font=\normalsize]
\foreach \x/\n/\c in {0/Arrendatario/Rama4, 5.2/App web (React)/AzulUNA, 10.4/API (Express)/Rama2, 15.6/Motor CASA SCORE/RaciR, 20.8/Base de datos/Rama5} {
  \node[draw=\c, fill=\c!15, line width=1.1pt, minimum width=4.2cm, minimum height=1.0cm, font=\large\bfseries] at (\x,0) {\n};
  \draw[dashed, gray, line width=0.8pt] (\x,-0.5) -- (\x,-12.2);}
\fill[Rama2!25] (10.25,-2.1) rectangle (10.55,-10.6);
\fill[RaciR!25] (15.45,-5.6) rectangle (15.75,-9.6);
\fill[Rama5!25] (20.65,-3.0) rectangle (20.95,-4.6);
\newcommand{\msg}[5]{\draw[->, line width=1.1pt, #5] (#1,#3) -- node[above, font=\normalsize]{#4} (#2,#3);}
\msg{0}{5.2}{-1.4}{1. elige pesos y punto}{}
\msg{5.2}{10.4}{-2.3}{2. GET /api/ranking?lat\&lng\&pesos}{}
\msg{10.4}{20.8}{-3.2}{3. leer cuartos activos}{}
\msg{20.8}{10.4}{-4.3}{4. cuartos, servicios, condiciones y fotos}{dashed}
\msg{10.4}{15.6}{-5.8}{5. calcular(cuartos, pesos, punto)}{}
\draw[line width=1pt, DoradoFIEI] (14.0,-6.4) rectangle (19.6,-8.9);
\node[draw=DoradoFIEI, fill=DoradoFIEI!20, font=\small\bfseries, anchor=north west] at (14.0,-6.4) {loop: cada cuarto};
\draw[->, line width=1.1pt] (15.75,-7.0) -- ++(1.2,0) -- ++(0,-0.6) node[right=4pt, font=\normalsize, align=left, midway, fill=white]{6. distancia, normalizar,\\bonos y penalidades} -- ++(-1.2,0);
\msg{15.6}{10.4}{-9.4}{7. lista ordenada con desglose}{dashed}
\msg{10.4}{5.2}{-10.4}{8. JSON del ranking}{dashed}
\msg{5.2}{0}{-11.4}{9. muestra la lista y los puntajes}{dashed}
\end{tikzpicture}
\end{grande}
\vspace{0.2cm}
{\large\noindent Flecha continua: llamada. Flecha discontinua: respuesta. El mismo motor se usa en la app Android, donde el cálculo se hace en el celular con los cuartos descargados.}
\end{hojah}

\begin{hojah}{UML: Despliegue y Estados del Cuarto}
\begin{grande}
\begin{tikzpicture}[>=Stealth, font=\normalsize, nodo/.style={draw=#1, line width=1.2pt, fill=#1!6, align=left, inner sep=8pt}, comp/.style={draw=#1, fill=white, line width=0.9pt, minimum width=3.6cm, minimum height=0.9cm, align=center}]
\node[font=\Large\bfseries, text=AzulUNA] at (5.5,7.6) {Despliegue};
\node[nodo=Rama4, minimum width=5.2cm, minimum height=3.4cm] (cel) at (0,4.6) {};
\node[font=\large\bfseries, anchor=north west, text=Rama4] at (cel.north west) {\guillemotleft dispositivo\guillemotright\ Celular Android};
\node[comp=Rama4] at (0,4.1) {App Expo + motor};
\node[comp=Rama4] at (0,3.1) {Caché SQLite};
\node[nodo=AzulUNA, minimum width=5.2cm, minimum height=3.4cm] (nav) at (0,0) {};
\node[font=\large\bfseries, anchor=north west, text=AzulUNA] at (nav.north west) {\guillemotleft dispositivo\guillemotright\ Navegador};
\node[comp=AzulUNA] at (0,-0.5) {Web React (Vite)};
\node[nodo=Rama2, minimum width=5.2cm, minimum height=3.4cm] (srv) at (7.0,0) {};
\node[font=\large\bfseries, anchor=north west, text=Rama2] at (srv.north west) {\guillemotleft servidor\guillemotright\ API};
\node[comp=Rama2] at (7.0,-0.1) {Express 5};
\node[comp=Rama2] at (7.0,-1.1) {Motor CASA SCORE};
\node[nodo=Rama5, minimum width=5.2cm, minimum height=3.4cm] (sb) at (7.0,4.6) {};
\node[font=\large\bfseries, anchor=north west, text=Rama5] at (sb.north west) {\guillemotleft nube\guillemotright\ Supabase};
\node[comp=Rama5] at (7.0,4.6) {Auth y PostgreSQL};
\node[comp=Rama5] at (7.0,3.6) {Storage (fotos)};
\node[nodo=GrisOscuro, minimum width=5.2cm, minimum height=1.6cm, align=center, font=\large\bfseries] (osm) at (3.5,-4.0) {OpenStreetMap\\{\normalsize\mdseries teselas del mapa}};
\draw[<->, line width=1.2pt] (cel.east) -- node[etq, above]{HTTPS} (sb.west);
\draw[<->, line width=1.2pt] (nav.east) -- node[etq, above]{HTTPS / JSON} (srv.west);
\draw[<->, line width=1.2pt] (srv.north) -- node[etq, right]{SQL} (sb.south);
\draw[->, line width=1.2pt] (nav.south) -- (osm.north west);
\draw[->, line width=1.2pt] (cel.south west) to[out=-120,in=150] (osm.west);
\draw[gray, line width=1pt] (11.6,7.8) -- (11.6,-5.2);
\node[font=\Large\bfseries, text=AzulUNA] at (17.6,7.6) {Estados de un cuarto};
\fill (13.4,4.6) circle (0.25);
\node[draw=AzulUNA, fill=AzulClaro, line width=1.1pt, rounded corners=10pt, minimum width=3.4cm, minimum height=1.2cm, font=\large\bfseries] (bor) at (16.6,4.6) {Borrador};
\node[draw=Rama4, fill=Rama4!15, line width=1.1pt, rounded corners=10pt, minimum width=3.4cm, minimum height=1.2cm, font=\large\bfseries] (act) at (16.6,0.4) {Activo};
\node[draw=GrisOscuro, fill=gray!15, line width=1.1pt, rounded corners=10pt, minimum width=3.4cm, minimum height=1.2cm, font=\large\bfseries] (ocu) at (21.4,0.4) {Oculto};
\draw[->, line width=1.2pt] (13.65,4.6) -- node[etq, above]{nuevo} (bor.west);
\draw[->, line width=1.2pt] (bor.south) -- node[etq, right, align=left]{publicar\\{[}datos y 4 fotos completos{]}} (act.north);
\draw[->, line width=1.2pt] (bor.east) to[out=20,in=-20, looseness=4] node[etq, right]{guardar borrador} (bor.north east);
\draw[->, line width=1.2pt] ([yshift=0.25cm]act.east) -- node[etq, above]{ocultar} ([yshift=0.25cm]ocu.west);
\draw[->, line width=1.2pt] ([yshift=-0.25cm]ocu.west) -- node[etq, below]{mostrar} ([yshift=-0.25cm]act.east);
\draw[->, line width=1.2pt] (act.west) to[out=200,in=160, looseness=4] node[etq, left]{editar} (act.south west);
\end{tikzpicture}
\end{grande}
\vspace{0.2cm}
{\large\noindent Solo los cuartos \textbf{activos} entran al ranking. La seguridad por fila permite que cada propietario edite y oculte únicamente sus propios cuartos.}
\end{hojah}

\begin{hojah}{Modelo Híbrido: Ciclo de Vida y Sprint Semanal}
\begin{grande}
\begin{tikzpicture}[>=Stealth, font=\normalsize]
% ciclo predictivo
\node[font=\Large\bfseries, text=AzulUNA, anchor=west] at (0,6.4) {Parte predictiva: fases con líneas base y gates};
\foreach \x/\w/\c/\t/\s in {0/4.6/Rama1/F1 Inicio y planificación/sem. 2 a 5, 4.8/4.6/Rama2/F2 MVP de la Unidad I/sem. 4 a 8, 9.6/4.6/Rama3/F3 Nube y nuevas funciones/sem. 10 a 13, 14.4/4.6/Rama4/F4 Piloto y aceptación/sem. 13 a 15, 19.2/4.6/Rama5/F5 Entrega y cierre/sem. 15 a 17} {
 \node[fill=\c, text=white, minimum width=\w cm, minimum height=1.3cm, align=center, anchor=west, font=\small\bfseries, text width=4.3cm] at (\x,5.0) {\t\\{\small\mdseries \s}};}
\foreach \x/\g/\d in {4.7/G1/11-09, 9.5/G2/09-10, 14.3/G3/27-11, 19.1/G4/04-12, 23.9/G5/18-12} {
 \node[diamond, fill=RaciR, inner sep=3pt] at (\x,3.9) {}; \node[font=\small\bfseries, text=RaciR, below] at (\x,3.7) {\g\ \d};}
% sprint loop
\node[font=\Large\bfseries, text=AzulUNA, anchor=west] at (0,2.2) {Parte ágil: sprint de una semana};
\node[circle, draw=AzulUNA, line width=1.4pt, fill=AzulClaro, minimum size=2.8cm, align=center, font=\large\bfseries] (inc) at (5.2,-3.0) {Incremento\\{\normalsize\mdseries probado}};
\foreach \a/\t/\c in {90/{Lunes\\planificación}/Rama1, 18/{Martes 20:00\\control: EVM y riesgos}/RaciR, -54/{Mié. y jue.\\desarrollo y documentos}/Rama3, -126/{Viernes 20:00\\revisión del sprint}/Rama4, 162/{Viernes\\retrospectiva}/Rama5} {
 \node[draw=\c, fill=\c!12, line width=1.1pt, rounded corners=5pt, align=center, font=\normalsize, minimum width=3.6cm, minimum height=1.2cm] at ($(inc)+(\a:4.1cm)$) {\t};}
\draw[->, line width=2pt, AzulUNA!60] ($(inc)+(80:2.6cm)$) arc[start angle=80, end angle=-260, radius=2.6cm];
% kanban
\node[font=\Large\bfseries, text=AzulUNA, anchor=west] at (11.2,1.6) {Tablero Kanban en GitHub Projects};
\foreach \i/\t/\s/\c in {0/Backlog/requisitos y tareas futuras/GrisOscuro, 1/Por hacer/comprometido en el sprint/Rama2, 2/En Progreso/límite: 2 tarjetas por integrante/SevMedia, 3/En Revisión/Pull Request o revisión cruzada/Rama5, 4/Hecho/criterio cumplido y evidencia/Rama4} {
 \node[draw=\c, fill=\c!12, line width=1.1pt, minimum width=2.45cm, minimum height=4.6cm, anchor=north] (k\i) at (12.5+\i*2.65,1.0) {};
 \node[fill=\c, text=white, font=\small\bfseries, minimum width=2.45cm, minimum height=0.7cm, anchor=north] at (12.5+\i*2.65,1.0) {\t};
 \node[font=\small, text width=2.2cm, align=center, anchor=north] at (12.5+\i*2.65,0.1) {\s};}
\foreach \i in {0,1,2,3} {\draw[->, line width=1.4pt, AzulUNA] (12.5+\i*2.65+1.25,-1.6) -- (12.5+\i*2.65+1.4,-1.6);}
\node[font=\normalsize, text width=12.8cm, align=left, anchor=north west] at (11.2,-3.9) {Cada tarjeta tiene responsable, apoyo, sprint, paquete de la EDT, estimación, horas reales y fecha límite. Pasa a Hecho solo con su enlace al commit o al Pull Request.};
\end{tikzpicture}
\end{grande}
\end{hojah}

\begin{hojah}{Comunicación con el Sponsor}
\begin{grande}
\begin{tikzpicture}[fl]
\node[actor=RaciR, minimum width=5.0cm, minimum height=1.8cm] (sp) at (0,0) {\textbf{Sponsor}\\{\normalsize Team Sobrevivientes}\\{\small decide por votación}};
\node[actor=AzulUNA, minimum width=4.4cm, minimum height=1.4cm] (pm) at (-10.5,0) {\textbf{Project Manager}\\{\small Clyde}};
\node[actor=Rama4, minimum width=4.4cm, minimum height=1.4cm] (int) at (-10.5,-4.2) {\textbf{Integrantes}\\{\small tablero y horas cada día}};
\node[actor=GrisOscuro, minimum width=4.4cm, minimum height=1.4cm] (as) at (-10.5,4.2) {\textbf{Asesor}\\{\small Ing. Renzo Apaza Cutipa}};
\draw[fle] (int) -- node[etq, right, align=left]{avance, horas\\y bloqueos} (pm);
\draw[fle] ([yshift=0.3cm]pm.east) -- node[etq, above, align=center]{martes: informe\\con semáforo} ([yshift=0.3cm]sp.west);
\draw[fle, RaciR] ([yshift=-0.3cm]sp.west) -- node[etq, below, align=center]{decisiones y acuerdos\\(acta en 24 h)} ([yshift=-0.3cm]pm.east);
\draw[<->, line width=1.1pt, GrisOscuro, >=Stealth] (as) -- node[etq, right, align=left]{entregables\\y observaciones} (pm);
\node[draw=Rama4, fill=Rama4!12, line width=1.1pt, rounded corners=4pt, text width=6.4cm, align=left, font=\normalsize] (v) at (8.6,3.6) {\textbf{\textcolor{Rama4}{VERDE} (0.95 o más)}\\Se informa en el control y queda en el acta.};
\node[draw=SevMedia, fill=SevMedia!12, line width=1.1pt, rounded corners=4pt, text width=6.4cm, align=left, font=\normalsize] (a) at (8.6,1.2) {\textbf{\textcolor{SevMedia}{ÁMBAR} (0.90 a 0.94)}\\El PM aplica una acción correctiva dentro del sprint e informa el viernes.};
\node[draw=RaciR, fill=RaciR!10, line width=1.1pt, rounded corners=4pt, text width=6.4cm, align=left, font=\normalsize] (r) at (8.6,-1.4) {\textbf{\textcolor{RaciR}{ROJO} (menos de 0.90)}\\Plan de recuperación o solicitud de cambio; análisis en 3 días y votación del sponsor.};
\node[draw=AzulUNA, fill=AzulClaro, line width=1.1pt, rounded corners=4pt, text width=6.4cm, align=left, font=\normalsize] (g) at (8.6,-4.0) {\textbf{GATE}\\El PM presenta el entregable con su criterio de aceptación; el sponsor vota si pasa a la siguiente fase.};
\foreach \n in {v,a,r,g} {\draw[->, line width=1.1pt, gray] (sp.east) -- (\n.west);}
\end{tikzpicture}
\end{grande}
\vspace{0.2cm}
{\footnotesize
\rowcolors{2}{white}{GrisFila}
\noindent\begin{tabularx}{\linewidth}{@{}P{5.6cm} P{3.4cm} P{4.2cm} L@{}}
\toprule
\rowcolor{AzulClaro}
\textbf{Qué recibe el sponsor} & \textbf{Cuándo} & \textbf{Canal} & \textbf{Para qué} \\
\midrule
Informe de valor ganado (SPI, CPI) y riesgos & Martes 20:00 & Google Meet y acta en GitHub & Decidir si hace falta una acción o un cambio \\
Demostración del incremento & Viernes 20:00 & Google Meet & Aceptar o rechazar lo construido en el sprint \\
Solicitud de cambio con análisis de impacto & Cuando se presenta & Issue con etiqueta \texttt{cambio} & Aprobar por unanimidad un cambio mayor \\
Entregable de cada gate & G1 a G5 & Reunión y repositorio & Votar el paso a la siguiente fase \\
\bottomrule
\end{tabularx}
}
\end{hojah}

\begin{hojah}{Interesados: Poder e Interés}
\begin{minipage}[c]{11.0cm}
\begin{adjustbox}{width=\linewidth}
\begin{tikzpicture}[font=\large]
\fill[CuadSatisf!15] (0,4) rectangle (4,8); \fill[CuadCerca!15] (4,4) rectangle (8,8); \fill[CuadMonit!12] (0,0) rectangle (4,4); \fill[CuadInfo!12] (4,0) rectangle (8,4);
\draw[line width=1pt] (0,0) rectangle (8,8); \draw[line width=1pt] (4,0) -- (4,8); \draw[line width=1pt] (0,4) -- (8,4);
\node[font=\normalsize\bfseries, text=CuadSatisf] at (2,7.6) {Mantener satisfecho}; \node[font=\normalsize\bfseries, text=CuadCerca] at (6,7.6) {Gestionar de cerca};
\node[font=\normalsize\bfseries, text=CuadMonit] at (2,0.4) {Monitorear}; \node[font=\normalsize\bfseries, text=CuadInfo] at (6,0.4) {Mantener informado};
\draw[->, line width=1.2pt, >=Stealth] (-0.4,0) -- (-0.4,8.3) node[above]{Poder}; \draw[->, line width=1.2pt, >=Stealth] (0,-0.4) -- (8.3,-0.4) node[right]{Interés};
\foreach \p/\i/\n in {5/5/I1, 4/5/I2, 4/4/I3, 3/4/I4, 4/3/I5, 4/2/I6, 2/3/I7, 3/1/I8} {
  \pgfmathsetmacro{\X}{ifthenelse(\i<=3,\i*1.25-0.5,5+(\i-4)*1.7)} \pgfmathsetmacro{\Y}{ifthenelse(\p<=3,\p*1.25-0.5,5+(\p-4)*1.7)}
  \fill[AzulUNA] (\X,\Y) circle (4pt); \node[anchor=south west, font=\large\bfseries] at (\X+0.05,\Y+0.05) {\n};}
\end{tikzpicture}
\end{adjustbox}
\end{minipage}\hfill
\begin{minipage}[c]{11.6cm}
{\normalsize
\rowcolors{2}{white}{GrisFila}
\noindent\begin{tabularx}{\linewidth}{@{}C{0.8cm} L C{1.2cm}@{}}
\toprule
\rowcolor{AzulClaro}
\textbf{ID} & \textbf{Interesado y estrategia} & \textbf{P, I} \\
\midrule
I1 & Team Sobrevivientes (sponsor): decide en cada gate & 5, 5 \\
I2 & Arrendatarios estudiantes: entrevistas, piloto y encuesta & 4, 5 \\
I3 & Propietarios: visitas, guía de publicación y WhatsApp & 4, 4 \\
I4 & Trabajadores y familias: avances por WhatsApp & 3, 4 \\
I5 & Docente asesor: entregables según la ruta & 4, 3 \\
I6 & Concursos y ferias: postulación con el MVP & 4, 2 \\
I7 & Anunciantes locales: lista de negocios del año 1 & 2, 3 \\
I8 & Proveedores de servicios: revisión mensual de límites & 3, 1 \\
\bottomrule
\end{tabularx}
}
\end{minipage}
\end{hojah}

\begin{hojah}{Estructura de Desglose del Trabajo (EDT)}
\begin{grande}
\begin{tikzpicture}[font=\normalsize]
\node[draw=AzulUNA, fill=AzulUNA, text=white, font=\Large\bfseries, minimum width=7cm, minimum height=1.2cm, rounded corners=4pt] (root) at (12.8,0) {CasaScore (AlquilaPuno) -- 512 h};
\foreach \i/\t/\h/\c in {0/1 Gestión del proyecto/168 h/Rama1, 1/2 Análisis y diseño/64 h/Rama2, 2/3 Construcción del software/160 h/Rama3, 3/4 Pruebas y validación/72 h/Rama4, 4/5 Cierre/48 h/Rama5} {
  \node[draw=\c, fill=\c, text=white, font=\normalsize\bfseries, text width=4.4cm, align=center, minimum height=1.2cm, rounded corners=3pt] (b\i) at (\i*5.4+2.0,-2.2) {\t\\\h};
  \draw[line width=1.2pt] (root.south) -- ++(0,-0.5) -| (b\i.north);}
\newcommand{\pq}[4]{\node[draw=#1, fill=#1!10, line width=0.9pt, text width=4.0cm, align=left, font=\normalsize, inner sep=4pt, minimum height=1.15cm, anchor=north west] (#2) at (#3) {#4}; \draw[#1, line width=0.9pt] ($(#3)+(-0.25,0.0)$) |- (#2.west);}
\foreach \k/\t in {1/1.1 Gobernanza del equipo,2/1.2 Viabilidad y Acta,3/1.3 Planificación base,4/1.4 Planes de soporte,5/1.5 Seguimiento y control,6/1.6 Entregas de la Unidad II} {\pq{Rama1}{a\k}{0.0+0.3,-2.5-\k*1.35}{\t}}
\foreach \k/\t in {1/2.1 Ficha y requerimientos,2/2.2 Modelo de datos y arquitectura,3/2.3 Diseño de pantallas} {\pq{Rama2}{c\k}{5.4+0.3,-2.5-\k*1.35}{\t}}
\foreach \k/\t in {1/3.1 Motor CASA SCORE,2/3.2 Publicación de cuartos,3/3.3 Ranking y detalle,4/3.4 Nube y registro,5/3.5 Punto de referencia,6/3.6 Perfil y mis cuartos} {\pq{Rama3}{d\k}{10.8+0.3,-2.5-\k*1.35}{\t}}
\foreach \k/\t in {1/4.1 Pruebas del motor,2/4.2 Levantamiento de cuartos,3/4.3 Piloto con clientes,4/4.4 Pruebas de aceptación} {\pq{Rama4}{e\k}{16.2+0.3,-2.5-\k*1.35}{\t}}
\foreach \k/\t in {1/5.1 Manuales,2/5.2 Informe final y exposición,3/5.3 Lecciones y acta de cierre} {\pq{Rama5}{f\k}{21.6+0.3,-2.5-\k*1.35}{\t}}
\end{tikzpicture}
\end{grande}
\vspace{0.2cm}
{\large\noindent 23 paquetes de 3 a 5 días. La suma de las ramas es igual a las 512 horas del equipo (regla del 100\,\%).}
\end{hojah}

\begin{hojah}{Diagrama de Gantt Detallado}
\begin{grande}[14.2cm]
\begin{tikzpicture}
\fill[AzulClaro] (0,0.15) rectangle (16.00,1.05);
\node[font=\small\bfseries, text=AzulUNA] at (0.50,0.78) {S2};
\node[font=\scriptsize, text=GrisOscuro] at (0.50,0.38) {31/08};
\node[font=\small\bfseries, text=AzulUNA] at (1.50,0.78) {S3};
\node[font=\scriptsize, text=GrisOscuro] at (1.50,0.38) {07/09};
\node[font=\small\bfseries, text=AzulUNA] at (2.50,0.78) {S4};
\node[font=\scriptsize, text=GrisOscuro] at (2.50,0.38) {14/09};
\node[font=\small\bfseries, text=AzulUNA] at (3.50,0.78) {S5};
\node[font=\scriptsize, text=GrisOscuro] at (3.50,0.38) {21/09};
\node[font=\small\bfseries, text=AzulUNA] at (4.50,0.78) {S6};
\node[font=\scriptsize, text=GrisOscuro] at (4.50,0.38) {28/09};
\node[font=\small\bfseries, text=AzulUNA] at (5.50,0.78) {S7};
\node[font=\scriptsize, text=GrisOscuro] at (5.50,0.38) {05/10};
\node[font=\small\bfseries, text=AzulUNA] at (6.50,0.78) {S8};
\node[font=\scriptsize, text=GrisOscuro] at (6.50,0.38) {12/10};
\node[font=\small\bfseries, text=AzulUNA] at (7.50,0.78) {S9};
\node[font=\scriptsize, text=GrisOscuro] at (7.50,0.38) {19/10};
\node[font=\small\bfseries, text=AzulUNA] at (8.50,0.78) {S10};
\node[font=\scriptsize, text=GrisOscuro] at (8.50,0.38) {26/10};
\node[font=\small\bfseries, text=AzulUNA] at (9.50,0.78) {S11};
\node[font=\scriptsize, text=GrisOscuro] at (9.50,0.38) {02/11};
\node[font=\small\bfseries, text=AzulUNA] at (10.50,0.78) {S12};
\node[font=\scriptsize, text=GrisOscuro] at (10.50,0.38) {09/11};
\node[font=\small\bfseries, text=AzulUNA] at (11.50,0.78) {S13};
\node[font=\scriptsize, text=GrisOscuro] at (11.50,0.38) {16/11};
\node[font=\small\bfseries, text=AzulUNA] at (12.50,0.78) {S14};
\node[font=\scriptsize, text=GrisOscuro] at (12.50,0.38) {23/11};
\node[font=\small\bfseries, text=AzulUNA] at (13.50,0.78) {S15};
\node[font=\scriptsize, text=GrisOscuro] at (13.50,0.38) {30/11};
\node[font=\small\bfseries, text=AzulUNA] at (14.50,0.78) {S16};
\node[font=\scriptsize, text=GrisOscuro] at (14.50,0.38) {07/12};
\node[font=\small\bfseries, text=AzulUNA] at (15.50,0.78) {S17};
\node[font=\scriptsize, text=GrisOscuro] at (15.50,0.38) {14/12};
\draw[gray!25] (0.00,0.15) -- (0.00,-13.70);
\draw[gray!25] (1.00,0.15) -- (1.00,-13.70);
\draw[gray!25] (2.00,0.15) -- (2.00,-13.70);
\draw[gray!25] (3.00,0.15) -- (3.00,-13.70);
\draw[gray!25] (4.00,0.15) -- (4.00,-13.70);
\draw[gray!25] (5.00,0.15) -- (5.00,-13.70);
\draw[gray!25] (6.00,0.15) -- (6.00,-13.70);
\draw[gray!25] (7.00,0.15) -- (7.00,-13.70);
\draw[gray!25] (8.00,0.15) -- (8.00,-13.70);
\draw[gray!25] (9.00,0.15) -- (9.00,-13.70);
\draw[gray!25] (10.00,0.15) -- (10.00,-13.70);
\draw[gray!25] (11.00,0.15) -- (11.00,-13.70);
\draw[gray!25] (12.00,0.15) -- (12.00,-13.70);
\draw[gray!25] (13.00,0.15) -- (13.00,-13.70);
\draw[gray!25] (14.00,0.15) -- (14.00,-13.70);
\draw[gray!25] (15.00,0.15) -- (15.00,-13.70);
\draw[gray!25] (16.00,0.15) -- (16.00,-13.70);
\fill[gray!18] (7.00,0.15) rectangle (8.00,-13.70);
\fill[gray!18] (15.00,0.15) rectangle (16.00,-13.70);
\node[font=\tiny, text=GrisOscuro, rotate=90] at (7.50,-12.50) {exámenes};
\node[font=\tiny, text=GrisOscuro, rotate=90] at (15.50,-12.50) {exámenes};
\node[anchor=east, font=\small\bfseries, text=AzulUNA] at (-1.25,0.6) {Paquete de trabajo};
\node[font=\small\bfseries, text=AzulUNA] at (-0.65,0.6) {Resp.};
\node[anchor=east, font=\normalsize\bfseries, text=Rama1] at (-1.25,-0.50) {1 Gestión del proyecto};
\node[anchor=east, font=\small] at (-1.25,-1.00) {1.1 Gobernanza del equipo};
\node[font=\small\bfseries, text=GrisOscuro] at (-0.65,-1.00) {VQ};
\fill[gray!55] (0.00,-0.90) rectangle (1.71,-0.80);
\fill[Rama4!25] (0.00,-1.20) rectangle (1.71,-0.94);
\fill[Rama4] (0.00,-1.20) rectangle (1.71,-0.94);
\draw[Rama4] (0.00,-1.20) rectangle (1.71,-0.94);
\node[anchor=west, font=\scriptsize\bfseries] at (1.76,-1.05) {100\,\%};
\node[anchor=east, font=\small] at (-1.25,-1.50) {1.2 Viabilidad y Acta de Constitución};
\node[font=\small\bfseries, text=GrisOscuro] at (-0.65,-1.50) {CP};
\fill[gray!55] (1.00,-1.40) rectangle (1.71,-1.30);
\fill[Rama4!25] (1.00,-1.70) rectangle (1.71,-1.44);
\fill[Rama4] (1.00,-1.70) rectangle (1.71,-1.44);
\draw[Rama4] (1.00,-1.70) rectangle (1.71,-1.44);
\node[anchor=west, font=\scriptsize\bfseries] at (1.76,-1.55) {100\,\%};
\node[anchor=east, font=\small] at (-1.25,-2.00) {1.3 Planificación base};
\node[font=\small\bfseries, text=GrisOscuro] at (-0.65,-2.00) {RA};
\fill[gray!55] (2.00,-1.90) rectangle (3.14,-1.80);
\fill[Rama4!25] (2.00,-2.20) rectangle (8.71,-1.94);
\fill[Rama4] (2.00,-2.20) rectangle (8.71,-1.94);
\draw[Rama4] (2.00,-2.20) rectangle (8.71,-1.94);
\node[anchor=west, font=\scriptsize\bfseries] at (8.76,-2.05) {100\,\%};
\node[anchor=east, font=\small] at (-1.25,-2.50) {1.4 Planes de soporte};
\node[font=\small\bfseries, text=GrisOscuro] at (-0.65,-2.50) {VQ};
\fill[gray!55] (3.00,-2.40) rectangle (4.14,-2.30);
\fill[Rama4!25] (3.00,-2.70) rectangle (8.71,-2.44);
\fill[Rama4] (3.00,-2.70) rectangle (8.71,-2.44);
\draw[Rama4] (3.00,-2.70) rectangle (8.71,-2.44);
\node[anchor=west, font=\scriptsize\bfseries] at (8.76,-2.55) {100\,\%};
\node[anchor=east, font=\small] at (-1.25,-3.00) {1.5 Seguimiento y control};
\node[font=\small\bfseries, text=GrisOscuro] at (-0.65,-3.00) {CP};
\fill[gray!55] (4.00,-2.90) rectangle (15.43,-2.80);
\fill[pattern=north east lines, pattern color=Rama1] (4.00,-3.20) rectangle (15.71,-2.94); \draw[Rama1] (4.00,-3.20) rectangle (15.71,-2.94);
\node[anchor=west, font=\scriptsize\bfseries] at (15.76,-3.07) {continuo};
\node[anchor=east, font=\small] at (-1.25,-3.50) {1.6 Entregas de la Unidad II};
\node[font=\small\bfseries, text=GrisOscuro] at (-0.65,-3.50) {CP};
\fill[gray!55] (6.14,-3.40) rectangle (8.14,-3.30);
\fill[Rama4!25] (6.14,-3.70) rectangle (8.71,-3.44);
\fill[Rama4] (6.14,-3.70) rectangle (8.71,-3.44);
\draw[Rama4] (6.14,-3.70) rectangle (8.71,-3.44);
\node[anchor=west, font=\scriptsize\bfseries] at (8.76,-3.55) {100\,\%};
\node[anchor=east, font=\normalsize\bfseries, text=Rama2] at (-1.25,-4.00) {2 Análisis y diseño};
\node[anchor=east, font=\small] at (-1.25,-4.50) {2.1 Ficha técnica y requerimientos};
\node[font=\small\bfseries, text=GrisOscuro] at (-0.65,-4.50) {CP};
\fill[gray!55] (0.00,-4.40) rectangle (0.71,-4.30);
\fill[Rama4!25] (0.00,-4.70) rectangle (0.71,-4.44);
\fill[Rama4] (0.00,-4.70) rectangle (0.71,-4.44);
\draw[Rama4] (0.00,-4.70) rectangle (0.71,-4.44);
\node[anchor=west, font=\scriptsize\bfseries] at (0.76,-4.55) {100\,\%};
\node[anchor=east, font=\small] at (-1.25,-5.00) {2.2 Modelo de datos y arquitectura};
\node[font=\small\bfseries, text=GrisOscuro] at (-0.65,-5.00) {ET};
\fill[gray!55] (2.00,-4.90) rectangle (3.14,-4.80);
\fill[Rama4!25] (2.00,-5.20) rectangle (6.71,-4.94);
\fill[Rama4] (2.00,-5.20) rectangle (6.71,-4.94);
\draw[Rama4] (2.00,-5.20) rectangle (6.71,-4.94);
\node[anchor=west, font=\scriptsize\bfseries] at (6.76,-5.05) {100\,\%};
\node[anchor=east, font=\small] at (-1.25,-5.50) {2.3 Diseño de pantallas};
\node[font=\small\bfseries, text=GrisOscuro] at (-0.65,-5.50) {ET};
\fill[gray!55] (2.00,-5.40) rectangle (3.14,-5.30);
\fill[Rama4!25] (2.00,-5.70) rectangle (6.71,-5.44);
\fill[Rama4] (2.00,-5.70) rectangle (6.71,-5.44);
\draw[Rama4] (2.00,-5.70) rectangle (6.71,-5.44);
\node[anchor=west, font=\scriptsize\bfseries] at (6.76,-5.55) {100\,\%};
\node[anchor=east, font=\normalsize\bfseries, text=Rama3] at (-1.25,-6.00) {3 Construcción del software};
\node[anchor=east, font=\small] at (-1.25,-6.50) {3.1 Motor CASA SCORE};
\node[font=\small\bfseries, text=GrisOscuro] at (-0.65,-6.50) {ET};
\fill[gray!55] (2.00,-6.40) rectangle (4.14,-6.30);
\fill[Rama4!25] (2.00,-6.70) rectangle (3.71,-6.44);
\fill[Rama4] (2.00,-6.70) rectangle (3.71,-6.44);
\draw[Rama4] (2.00,-6.70) rectangle (3.71,-6.44);
\node[anchor=west, font=\scriptsize\bfseries] at (4.19,-6.55) {100\,\%};
\node[anchor=east, font=\small] at (-1.25,-7.00) {3.2 Publicación de cuartos};
\node[font=\small\bfseries, text=GrisOscuro] at (-0.65,-7.00) {ET};
\fill[gray!55] (3.00,-6.90) rectangle (4.14,-6.80);
\fill[Rama4!25] (3.00,-7.20) rectangle (5.71,-6.94);
\fill[Rama4] (3.00,-7.20) rectangle (5.71,-6.94);
\draw[Rama4] (3.00,-7.20) rectangle (5.71,-6.94);
\node[anchor=west, font=\scriptsize\bfseries] at (5.76,-7.05) {100\,\%};
\node[anchor=east, font=\small] at (-1.25,-7.50) {3.3 Ranking y detalle (MVP Unidad I)};
\node[font=\small\bfseries, text=GrisOscuro] at (-0.65,-7.50) {ET};
\fill[gray!55] (4.00,-7.40) rectangle (6.14,-7.30);
\fill[Rama4!25] (4.00,-7.70) rectangle (5.71,-7.44);
\fill[Rama4] (4.00,-7.70) rectangle (5.71,-7.44);
\draw[Rama4] (4.00,-7.70) rectangle (5.71,-7.44);
\node[anchor=west, font=\scriptsize\bfseries] at (6.19,-7.55) {100\,\%};
\node[anchor=east, font=\small] at (-1.25,-8.00) {3.4 Nube, registro e inicio de sesión};
\node[font=\small\bfseries, text=GrisOscuro] at (-0.65,-8.00) {ET};
\fill[gray!55] (6.14,-7.90) rectangle (10.29,-7.80);
\fill[SevMedia!25] (6.14,-8.20) rectangle (11.43,-7.94);
\fill[SevMedia] (6.14,-8.20) rectangle (8.26,-7.94);
\draw[SevMedia] (6.14,-8.20) rectangle (11.43,-7.94);
\node[anchor=west, font=\scriptsize\bfseries] at (11.48,-8.05) {40\,\%};
\node[anchor=east, font=\small] at (-1.25,-8.50) {3.5 Punto de referencia y criterios};
\node[font=\small\bfseries, text=GrisOscuro] at (-0.65,-8.50) {ET};
\fill[gray!55] (6.14,-8.40) rectangle (9.29,-8.30);
\fill[SevMedia!25] (6.14,-8.70) rectangle (10.71,-8.44);
\fill[SevMedia] (6.14,-8.70) rectangle (7.06,-8.44);
\draw[SevMedia] (6.14,-8.70) rectangle (10.71,-8.44);
\node[anchor=west, font=\scriptsize\bfseries] at (10.76,-8.55) {20\,\%};
\node[anchor=east, font=\small] at (-1.25,-9.00) {3.6 Perfil y mis cuartos};
\node[font=\small\bfseries, text=GrisOscuro] at (-0.65,-9.00) {ET};
\fill[gray!55] (10.14,-8.90) rectangle (11.29,-8.80);
\fill[Rama2!25] (10.71,-9.20) rectangle (11.71,-8.94);
\draw[Rama2] (10.71,-9.20) rectangle (11.71,-8.94);
\node[anchor=west, font=\scriptsize\bfseries] at (11.76,-9.05) {0\,\%};
\node[anchor=east, font=\normalsize\bfseries, text=Rama4] at (-1.25,-9.50) {4 Pruebas y validación};
\node[anchor=east, font=\small] at (-1.25,-10.00) {4.1 Pruebas del motor};
\node[font=\small\bfseries, text=GrisOscuro] at (-0.65,-10.00) {VQ};
\fill[gray!55] (4.00,-9.90) rectangle (5.14,-9.80);
\fill[Rama4!25] (4.00,-10.20) rectangle (5.71,-9.94);
\fill[Rama4] (4.00,-10.20) rectangle (5.71,-9.94);
\draw[Rama4] (4.00,-10.20) rectangle (5.71,-9.94);
\node[anchor=west, font=\scriptsize\bfseries] at (5.76,-10.05) {100\,\%};
\node[anchor=east, font=\small] at (-1.25,-10.50) {4.2 Levantamiento de cuartos};
\node[font=\small\bfseries, text=GrisOscuro] at (-0.65,-10.50) {CP};
\fill[gray!55] (3.00,-10.40) rectangle (4.14,-10.30);
\fill[SevMedia!25] (3.00,-10.70) rectangle (11.86,-10.44);
\fill[SevMedia] (3.00,-10.70) rectangle (8.93,-10.44);
\draw[SevMedia] (3.00,-10.70) rectangle (11.86,-10.44);
\node[anchor=west, font=\scriptsize\bfseries] at (11.91,-10.55) {67\,\%};
\node[anchor=east, font=\small] at (-1.25,-11.00) {4.3 Piloto con clientes};
\node[font=\small\bfseries, text=GrisOscuro] at (-0.65,-11.00) {CP};
\fill[gray!55] (10.14,-10.90) rectangle (11.29,-10.80);
\fill[Rama2!25] (12.00,-11.20) rectangle (12.71,-10.94);
\draw[Rama2] (12.00,-11.20) rectangle (12.71,-10.94);
\node[anchor=west, font=\scriptsize\bfseries] at (12.76,-11.05) {0\,\%};
\node[anchor=east, font=\small] at (-1.25,-11.50) {4.4 Pruebas de aceptación};
\node[font=\small\bfseries, text=GrisOscuro] at (-0.65,-11.50) {ET};
\fill[gray!55] (11.29,-11.40) rectangle (12.43,-11.30);
\fill[Rama2!25] (13.00,-11.70) rectangle (13.71,-11.44);
\draw[Rama2] (13.00,-11.70) rectangle (13.71,-11.44);
\node[anchor=west, font=\scriptsize\bfseries] at (13.76,-11.55) {0\,\%};
\node[anchor=east, font=\normalsize\bfseries, text=Rama5] at (-1.25,-12.00) {5 Cierre};
\node[anchor=east, font=\small] at (-1.25,-12.50) {5.1 Manuales};
\node[font=\small\bfseries, text=GrisOscuro] at (-0.65,-12.50) {ET};
\fill[gray!55] (12.29,-12.40) rectangle (13.43,-12.30);
\fill[Rama2!25] (13.00,-12.70) rectangle (13.71,-12.44);
\draw[Rama2] (13.00,-12.70) rectangle (13.71,-12.44);
\node[anchor=west, font=\scriptsize\bfseries] at (13.76,-12.55) {0\,\%};
\node[anchor=east, font=\small] at (-1.25,-13.00) {5.2 Informe final y exposición};
\node[font=\small\bfseries, text=GrisOscuro] at (-0.65,-13.00) {CP};
\fill[gray!55] (13.29,-12.90) rectangle (15.14,-12.80);
\fill[Rama2!25] (13.29,-13.20) rectangle (15.14,-12.94);
\draw[Rama2] (13.29,-13.20) rectangle (15.14,-12.94);
\node[anchor=west, font=\scriptsize\bfseries] at (15.19,-13.05) {0\,\%};
\node[anchor=east, font=\small] at (-1.25,-13.50) {5.3 Lecciones aprendidas y acta de cierre};
\node[font=\small\bfseries, text=GrisOscuro] at (-0.65,-13.50) {VQ};
\fill[gray!55] (15.00,-13.40) rectangle (15.43,-13.30);
\fill[Rama2!25] (15.00,-13.70) rectangle (15.71,-13.44);
\draw[Rama2] (15.00,-13.70) rectangle (15.71,-13.44);
\node[anchor=west, font=\scriptsize\bfseries] at (15.76,-13.55) {0\,\%};
\node[anchor=east, font=\normalsize\bfseries, text=RaciR] at (-1.25,-14.05) {Hitos};
\node[diamond, fill=RaciR, inner sep=2.5pt] at (1.71,-14.05) {};
\node[font=\scriptsize\bfseries, text=RaciR, above=1pt] at (1.71,-13.93) {H1};
\node[diamond, fill=RaciR, inner sep=2.5pt] at (3.29,-14.05) {};
\node[font=\scriptsize\bfseries, text=RaciR, above=1pt] at (3.29,-13.93) {H2};
\node[diamond, fill=RaciR, inner sep=2.5pt] at (5.71,-14.05) {};
\node[font=\scriptsize\bfseries, text=RaciR, above=1pt] at (5.71,-13.93) {H3};
\node[diamond, fill=RaciR, inner sep=2.5pt] at (8.14,-14.05) {};
\node[font=\scriptsize\bfseries, text=RaciR, above=1pt] at (8.14,-13.93) {H4};
\node[diamond, draw=gray, inner sep=2pt] at (10.29,-14.05) {};
\node[diamond, fill=RaciR, inner sep=2.5pt] at (11.43,-14.05) {};
\node[font=\scriptsize\bfseries, text=RaciR, above=1pt] at (11.43,-13.93) {H5};
\node[diamond, draw=gray, inner sep=2pt] at (11.29,-14.05) {};
\node[diamond, fill=RaciR, inner sep=2.5pt] at (12.71,-14.05) {};
\node[font=\scriptsize\bfseries, text=RaciR, above=1pt] at (12.71,-13.93) {H6};
\node[diamond, draw=gray, inner sep=2pt] at (12.43,-14.05) {};
\node[diamond, fill=RaciR, inner sep=2.5pt] at (13.71,-14.05) {};
\node[font=\scriptsize\bfseries, text=RaciR, above=1pt] at (13.71,-13.93) {H7};
\node[diamond, fill=RaciR, inner sep=2.5pt] at (15.14,-14.05) {};
\node[font=\scriptsize\bfseries, text=RaciR, above=1pt] at (15.14,-13.93) {H8};
\node[diamond, draw=gray, inner sep=2pt] at (15.43,-14.05) {};
\node[diamond, fill=RaciR, inner sep=2.5pt] at (15.71,-14.05) {};
\node[font=\scriptsize\bfseries, text=RaciR, above=1pt] at (15.71,-13.93) {H9};
\draw[RaciR, line width=1.3pt, dashed] (10.00,1.1) -- (10.00,-14.40) node[below, font=\small\bfseries, text=RaciR]{hoy 09/11};
\end{tikzpicture}
\end{grande}
\vspace{0.1cm}
{\footnotesize\noindent \tikz{\fill[gray!55] (0,0) rectangle (0.5,0.1);}~Línea base \quad \tikz{\fill[Rama4] (0,0) rectangle (0.5,0.25);}~Terminado \quad \tikz{\fill[SevMedia] (0,0) rectangle (0.25,0.25); \fill[SevMedia!25] (0.25,0) rectangle (0.5,0.25);}~En curso (parte llena = avance) \quad \tikz{\fill[Rama2!25] (0,0) rectangle (0.5,0.25);}~Pendiente (fechas reprogramadas) \quad \tikz{\node[diamond, fill=RaciR, inner sep=2pt]{};}~Hito \quad \tikz{\node[diamond, draw=gray, inner sep=2pt]{};}~Hito en línea base \quad Responsables: CP Clyde, RA Ruth, ET Eliseo, VQ Viviana. Avance al corte del 06/11.}
\end{hojah}

\begin{hojah}{Red de Actividades y Ruta Crítica}
\begin{grande}
\begin{tikzpicture}[>=Stealth, nd/.style={draw=#1, fill=#1!10, line width=1pt, rounded corners=3pt, minimum width=2.0cm, minimum height=1.2cm, align=center, font=\normalsize, inner sep=2pt}]
\newcommand{\nd}[6]{\node[nd=#5, #6] (#1) at (#3) {\textbf{#2}\\{\small #4}};}
\nd{n21}{2.1}{0,0}{1.00 sem}{RaciR}{line width=2pt}
\nd{n11}{1.1}{0,-2.2}{2.00 sem}{Rama1}{}
\nd{n12}{1.2}{2.8,0}{1.00 sem}{RaciR}{line width=2pt}
\nd{n22}{2.2}{5.6,1.6}{1.05 sem}{RaciR}{line width=2pt}
\nd{n23}{2.3}{5.6,0}{1.05 sem}{RaciR}{line width=2pt}
\nd{n31}{3.1}{5.6,-1.6}{2.10 sem}{RaciR}{line width=2pt}
\nd{n13}{1.3}{5.6,-3.6}{1.05 sem}{Rama1}{}
\nd{n32}{3.2}{8.4,0.8}{1.05 sem}{RaciR}{line width=2pt}
\nd{n42}{4.2}{8.4,-1.0}{1.05 sem}{Rama4}{}
\nd{n14}{1.4}{8.4,-3.6}{1.05 sem}{Rama1}{}
\nd{n33}{3.3}{11.2,0.8}{2.10 sem}{RaciR}{line width=2pt}
\nd{n41}{4.1}{11.2,-1.0}{1.05 sem}{Rama4}{}
\nd{n34}{3.4}{14.0,1.6}{3.15 sem}{RaciR}{line width=2pt}
\nd{n35}{3.5}{14.0,-0.4}{2.10 sem}{Rama3}{}
\nd{n16}{1.6}{14.0,-2.4}{1.89 sem}{Rama1}{}
\nd{n36}{3.6}{16.8,1.6}{1.05 sem}{RaciR}{line width=2pt}
\nd{n43}{4.3}{16.8,-0.4}{1.05 sem}{RaciR}{line width=2pt}
\nd{n44}{4.4}{19.6,0.6}{1.05 sem}{RaciR}{line width=2pt}
\nd{n51}{5.1}{19.6,-1.6}{1.05 sem}{RaciR}{line width=2pt}
\nd{n52}{5.2}{22.4,-1.6}{1.57 sem}{RaciR}{line width=2pt}
\nd{n53}{5.3}{22.4,-3.6}{0.53 sem}{RaciR}{line width=2pt}
\foreach \a/\b in {n21/n12, n12/n22, n12/n23, n12/n31, n22/n32, n23/n32, n31/n33, n32/n33, n33/n34, n34/n36, n34/n43, n36/n44, n43/n44, n44/n51, n51/n52, n52/n53} {\draw[->, line width=1.6pt, RaciR] (\a) -- (\b);}
\foreach \a/\b in {n12/n13, n13/n14, n23/n42, n31/n41, n32/n41, n33/n35, n35/n36, n35/n43, n42/n43, n41/n44, n33/n16} {\draw[->, line width=1pt, gray] (\a) -- (\b);}
\end{tikzpicture}
\end{grande}
\vspace{0.3cm}
{\large\noindent \textcolor{RaciR}{\textbf{Ruta crítica}} (holgura cero): 2.1, 1.2, 2.2, 2.3, 3.1, 3.2, 3.3, 3.4, 3.6, 4.3, 4.4, 5.1, 5.2 y 5.3. Dura 15.60 semanas y termina el 16/12/2026. Con $t_e = (o + 4m + p)/6$ y $\sigma = 0.60$ semanas, la probabilidad de terminar hasta el 18/12 es del \textbf{75\,\%}. Hoy la actividad crítica atrasada es la 3.4 (nube y registro).}
\end{hojah}

\begin{hojah}{Valor Ganado: Curva S al Corte de la Semana 11}
\begin{grande}
\begin{tikzpicture}[font=\large, >=Stealth]
\foreach \v/\l in {0/0,2/2\,000,4/4\,000,6/6\,000,8/8\,000,10/10\,000,12/12\,000} {\draw[gray!25] (0,\v) -- (20,\v); \node[anchor=east, font=\normalsize] at (-0.1,\v) {\l};}
\foreach \w in {1,...,17} {\node[font=\normalsize] at ({(\w-1)*1.25},-0.45) {\w};}
\node[font=\normalsize\bfseries] at (10,-1.05) {Semana del semestre};
\node[rotate=90, font=\normalsize\bfseries] at (-1.7,6) {Soles acumulados};
\draw[line width=0.8pt] (0,0) -- (20.3,0); \draw[line width=0.8pt] (0,0) -- (0,12.4);
\draw[GrisOscuro, dashed, line width=1pt] (0,11.108) -- (20,11.108) node[pos=0.02, above right, font=\normalsize\bfseries]{BAC del trabajo: S/ 11,108.07};
\draw[RaciR, dashed, line width=1pt] (6.25,0) -- (6.25,12.2) node[above, font=\normalsize\bfseries, text=RaciR]{corte 1 (02/10)};
\draw[RaciR, dashed, line width=1pt] (12.5,0) -- (12.5,12.2) node[above, font=\normalsize\bfseries, text=RaciR]{corte 2 (06/11)};
\draw[AzulUNA, line width=2.6pt] (0.00,0.000) -- (1.25,0.559) -- (2.50,1.117) -- (3.75,2.420) -- (5.00,3.812) -- (6.25,4.691) -- (7.50,5.313) -- (8.75,6.061) -- (10.00,6.349) -- (11.25,6.960) -- (12.50,7.385) -- (13.75,8.230) -- (15.00,9.061) -- (16.25,9.530) -- (17.50,10.089) -- (18.75,10.876) -- (20.00,11.108);
\draw[Rama4, line width=2.6pt] (0.00,0.000) -- (1.25,0.559) -- (2.50,1.117) -- (3.75,1.649) -- (5.00,2.149) -- (6.25,2.625) -- (7.50,3.611) -- (8.75,4.521) -- (10.00,5.149) -- (11.25,5.860) -- (12.50,6.426);
\draw[RaciR, line width=2.6pt, dashed] (0.00,0.000) -- (1.25,0.559) -- (2.50,1.126) -- (3.75,1.704) -- (5.00,2.247) -- (6.25,2.736) -- (7.50,3.794) -- (8.75,4.761) -- (10.00,5.431) -- (11.25,6.170) -- (12.50,6.769);
\draw[RaciR!60, line width=1.6pt, dotted] (12.5,6.769) -- (20,11.700) node[right, font=\normalsize\bfseries, text=RaciR]{EAC 11,700.40};
\node[right, font=\normalsize\bfseries, text=AzulUNA] at (12.9,7.55) {PV 7,385.01};
\node[right, font=\normalsize\bfseries, text=RaciR] at (12.9,6.85) {AC 6,768.95};
\node[right, font=\normalsize\bfseries, text=Rama4] at (12.9,6.15) {EV 6,426.28};
\draw[<->, line width=1.3pt] (12.25,6.426) -- node[left=8pt, font=\normalsize\bfseries, fill=white]{SV $-$958.73} (12.25,7.385);
\end{tikzpicture}
\end{grande}
\vspace{0.1cm}
{\large\noindent \textcolor{AzulUNA}{\rule{18pt}{3pt}}~PV: lo que debía estar hecho \quad \textcolor{Rama4}{\rule{18pt}{3pt}}~EV: lo realmente hecho \quad \textcolor{RaciR}{\rule{18pt}{3pt}}~AC: lo que costó. Entre los dos cortes el SPI subió de 0.56 a 0.87 y el CPI se mantuvo entre 0.96 y 0.95.}
\end{hojah}

\begin{hojah}{Tablero de Indicadores al 06/11/2026}
\begin{grande}
\begin{tikzpicture}
\newcommand{\kpi}[6]{\node[draw=#3, fill=#3!9, line width=1.4pt, rounded corners=5pt, minimum width=5.4cm, minimum height=3.0cm] at (#1,#2) {};
\node[font=\small\bfseries, text=GrisOscuro] at (#1,#2+0.95) {#4};
\node[font=\Huge\bfseries, text=#3] at (#1,#2+0.05) {#5};
\node[font=\small, text=GrisOscuro, text width=5.0cm, align=center] at (#1,#2-0.95) {#6};}
\node[font=\Large\bfseries, text=AzulUNA, anchor=west] at (-2.7,2.2) {Cronograma y costo};
\kpi{0}{0}{RaciR}{SPI = EV / PV}{0.87}{rojo; era 0.56 en la semana 6}
\kpi{5.8}{0}{Rama4}{CPI = EV / AC}{0.95}{verde en el límite}
\kpi{11.6}{0}{RaciR}{SV = EV $-$ PV}{$-$958.73}{atraso de cerca de una semana}
\kpi{17.4}{0}{SevMedia}{CV = EV $-$ AC}{$-$342.68}{ligero sobrecosto}
\kpi{0}{-3.4}{SevMedia}{EAC = BAC / CPI}{11,700.40}{costo final pronosticado}
\kpi{5.8}{-3.4}{SevMedia}{TCPI}{1.08}{eficiencia necesaria desde hoy}
\kpi{11.6}{-3.4}{SevMedia}{AVANCE REAL}{57.9\,\%}{planificado: 66.5\,\%}
\kpi{17.4}{-3.4}{Rama4}{CONTINGENCIA USADA}{S/ 0}{de S/ 555.40 disponibles}
\node[font=\Large\bfseries, text=AzulUNA, anchor=west] at (-2.7,-5.6) {Producto, riesgos y cambios};
\kpi{0}{-7.8}{Rama4}{PRUEBAS DEL MOTOR}{7 / 7}{aprobadas; ranking en menos de 250 ms}
\kpi{5.8}{-7.8}{SevMedia}{REQUISITOS TERMINADOS}{6 / 9}{faltan registro, perfil y nube}
\kpi{11.6}{-7.8}{SevMedia}{CUARTOS PARA EL PILOTO}{10 / 12}{meta propuesta de 12}
\kpi{17.4}{-7.8}{Rama4}{INSCRITOS AL PILOTO}{14 / 10}{meta cubierta}
\kpi{0}{-11.2}{RaciR}{RIESGOS ALTOS ABIERTOS}{3}{R02 activado (9), R01 y R08 (6)}
\kpi{5.8}{-11.2}{SevMedia}{CAMBIOS PENDIENTES}{1}{SC-02: piloto a la semana 14}
\kpi{11.6}{-11.2}{Rama4}{HITOS CUMPLIDOS}{4 / 9}{H1 a H4 a tiempo; H5 al 18/11}
\kpi{17.4}{-11.2}{Rama4}{DOCUMENTOS EN GITHUB}{21 / 21}{todos los previstos hasta hoy}
\end{tikzpicture}
\end{grande}
\vspace{0.1cm}
{\large\noindent \textbf{Semáforo:} \etiqueta{Rama4}{VERDE} 0.95 o más o meta cumplida \quad \etiqueta{SevMedia}{ÁMBAR} de 0.90 a 0.94 o meta en riesgo \quad \etiqueta{RaciR}{ROJO} menos de 0.90 o meta incumplida. Los indicadores se calculan cada martes con la vista Control semanal del tablero.}
\end{hojah}

\begin{hojah}{Presupuesto: de la Estimación al Presupuesto Total}
\begin{grande}
\begin{tikzpicture}[font=\large]
\def\s{0.00085}
\foreach \v/\l in {0/0,2/2\,000,4/4\,000,6/6\,000,8/8\,000,10/10\,000,12/12\,000} {\draw[gray!25] (0,\v*1000*\s) -- (23,\v*1000*\s); \node[anchor=east, font=\normalsize] at (-0.1,\v*1000*\s) {\l};}
\draw[line width=0.8pt] (0,0) -- (23,0);
\newcommand{\barra}[6]{\fill[#5] (#1,#2*\s) rectangle (#1+2.6,#3*\s); \node[font=\normalsize\bfseries, above] at (#1+1.3,#3*\s) {#4}; \node[font=\normalsize, text width=3.0cm, align=center, below] at (#1+1.3,-0.1) {#6};}
\barra{0.4}{0}{7681.40}{7,681.40}{Rama2}{Personal: 512 h con EsSalud}
\barra{3.6}{7681.40}{8600.63}{919.23}{Rama3}{Depreciación de equipos}
\barra{6.8}{8600.63}{11108.07}{2,507.44}{DoradoFIEI}{Gasto en efectivo}
\barra{10.0}{0}{11108.07}{11,108.07}{AzulUNA}{Estimación de costos}
\barra{13.2}{11108.07}{11663.47}{555.40}{SevMedia}{Contingencia (7 riesgos)}
\barra{16.4}{0}{11663.47}{11,663.47}{AzulUNA}{Línea base (BAC)}
\barra{19.6}{11663.47}{12829.82}{1,166.35}{RaciR}{Reserva de gestión (10\,\%)}
\draw[dashed, gray] (3.0,7681.40*\s) -- (3.6,7681.40*\s); \draw[dashed, gray] (6.2,8600.63*\s) -- (6.8,8600.63*\s);
\draw[dashed, gray] (12.6,11108.07*\s) -- (13.2,11108.07*\s); \draw[dashed, gray] (19.0,11663.47*\s) -- (19.6,11663.47*\s);
\node[draw=RaciR, line width=1.3pt, fill=white, font=\large\bfseries, align=center] at (21.0,6.5) {Presupuesto total\\S/ 12,829.82};
\end{tikzpicture}
\end{grande}
\vspace{0.2cm}
{\large\noindent La contingencia cubre los riesgos conocidos y la usa el Project Manager; la reserva de gestión cubre lo imprevisto y solo se libera con un cambio aprobado por el sponsor. De los S/ 12,829.82, solo S/ 2,507.44 salen del bolsillo del equipo; el resto valoriza el tiempo y los equipos.}
\end{hojah}

\begin{hojah}{Ingresos de la App: Modelo de Negocio}
\begin{minipage}[c]{12.4cm}
\begin{adjustbox}{width=\linewidth}
\begin{tikzpicture}[font=\large]
\foreach \v/\l in {0/0,2/2\,000,4/4\,000,6/6\,000,8/8\,000,10/10\,000} {\draw[gray!25] (0,\v*1000*0.00105) -- (13.5,\v*1000*0.00105); \node[anchor=east, font=\normalsize] at (-0.1,\v*1000*0.00105) {\l};}
\draw[line width=1pt] (0,0) -- (13.5,0);
\fill[Rama2] (1.00,0.0000) rectangle (3.40,1.2049);
\fill[DoradoFIEI] (1.00,1.2049) rectangle (3.40,2.1499);
\fill[RaciR] (1.00,2.1499) rectangle (3.40,2.4746);
\fill[Rama4] (1.00,2.4746) rectangle (3.40,2.5525);
\fill[Rama5] (1.00,2.5525) rectangle (3.40,2.5962);
\node[font=\large\bfseries, text=white, below] at (2.20,2.5962) {2,472.61};
\draw[RaciR, line width=2pt, dashed] (0.60,5.2845) -- (3.80,5.2845);
\node[font=\normalsize\bfseries, text=RaciR, above left] at (3.80,5.2845) {costo 5,033};
\node[font=\large\bfseries, below] at (2.20,-0.1) {Año 1 (2027)};
\fill[Rama2] (5.30,0.0000) rectangle (7.70,2.4097);
\fill[DoradoFIEI] (5.30,2.4097) rectangle (7.70,4.2997);
\fill[RaciR] (5.30,4.2997) rectangle (7.70,4.9492);
\fill[Rama4] (5.30,4.9492) rectangle (7.70,5.1050);
\fill[Rama5] (5.30,5.1050) rectangle (7.70,5.1925);
\node[font=\large\bfseries, text=white, below] at (6.50,5.1925) {4,945.22};
\draw[RaciR, line width=2pt, dashed] (4.90,7.5214) -- (8.10,7.5214);
\node[font=\normalsize\bfseries, text=RaciR, above left] at (8.10,7.5214) {costo 7,163};
\node[font=\large\bfseries, below] at (6.50,-0.1) {Año 2 (2028)};
\fill[Rama2] (9.60,0.0000) rectangle (12.00,4.8195);
\fill[DoradoFIEI] (9.60,4.8195) rectangle (12.00,8.5995);
\fill[RaciR] (9.60,8.5995) rectangle (12.00,9.8983);
\fill[Rama4] (9.60,9.8983) rectangle (12.00,10.2100);
\fill[Rama5] (9.60,10.2100) rectangle (12.00,10.3850);
\node[font=\large\bfseries, text=white, below] at (10.80,10.3850) {9,890.43};
\draw[RaciR, line width=2pt, dashed] (9.20,10.5062) -- (12.40,10.5062);
\node[font=\normalsize\bfseries, text=RaciR, above left] at (12.40,10.5062) {costo 10,006};
\node[font=\large\bfseries, below] at (10.80,-0.1) {Año 3 (2029)};
\end{tikzpicture}
\end{adjustbox}
{\small \textcolor{Rama2}{\rule{10pt}{10pt}}~Suscripción por cuarto \quad \textcolor{DoradoFIEI}{\rule{10pt}{10pt}}~Anunciantes locales \quad \textcolor{RaciR}{\rule{10pt}{10pt}}~Publicidad en internet\\ \textcolor{Rama4}{\rule{10pt}{10pt}}~Anuncios recompensados \quad \textcolor{Rama5}{\rule{10pt}{10pt}}~Versión Premium \quad \textcolor{RaciR}{\rule[2pt]{14pt}{1.5pt}}~Costo de operación}
\end{minipage}\hfill
\begin{minipage}[c]{10.8cm}
{\small
\setlength{\tabcolsep}{3pt}
\rowcolors{2}{white}{GrisFila}
\noindent\begin{tabularx}{\linewidth}{@{}P{2.6cm} L Q{1.25cm}@{}}
\toprule
\rowcolor{AzulClaro}
\textbf{Fuente} & \textbf{Supuesto} & \textbf{Año 1} \\
\midrule
Suscripción por cuarto & S/ 5 al mes por cuarto; 270 cuartos-mes & 1,147.50 \\
Anunciantes locales & 5 negocios de Puno $\times$ S/ 15 al mes & 900.00 \\
\textbf{Publicidad en internet} & \textbf{Banners de una red de anuncios en la web y la app: 30 páginas por usuario al mes, US\$ 0.50 por 1,000 impresiones} & \textbf{309.24} \\
Anuncios recompensados & 2 videos por usuario al mes; US\$ 1.8 por 1,000 & 74.22 \\
Versión Premium & 2\,\% de usuarios $\times$ S/ 4.90 & 41.65 \\
\midrule
\rowcolor{white}\textbf{Ingresos} & 500, 1,000 y 2,000 usuarios al mes & \textbf{2,472.61} \\
\rowcolor{white}\textbf{Costos de operación} & Mantenimiento, moderación, base de datos y difusión & \textbf{5,032.83} \\
\bottomrule
\end{tabularx}
}

\vspace{0.25cm}
{\normalsize\noindent La publicidad en internet paga por impresiones y no cambia el ranking: los banners van en espacios fijos, separados de los cuartos. Aun así, en los tres primeros años los ingresos no cubren la operación, por eso el reto es mantener el producto sin gastar más de lo que genera mientras crecen los usuarios.}
\end{minipage}
\end{hojah}

\begin{hojah}{Rentabilidad y Payback}
\begin{minipage}[c]{13.6cm}
\begin{adjustbox}{width=\linewidth}
\begin{tikzpicture}[x=4.2cm, y=0.075cm, font=\large]
\foreach \v in {-20,0,20,40,60,80} {\draw[gray!25] (0,\v) -- (3,\v); \node[anchor=east, font=\normalsize] at (-0.03,\v) {\v};}
\draw[line width=1.2pt] (0,0) -- (3.05,0);
\foreach \a in {0,1,2,3} {\node[font=\normalsize\bfseries] at (\a,-26) {Año \a};}
\node[rotate=90, font=\normalsize\bfseries] at (-0.33,30) {Miles de soles acumulados};
\draw[Rama4, line width=2.6pt] (0,-11.108) -- (1,0.906) -- (2,27.836) -- (3,86.016);
\draw[AzulUNA, line width=2.6pt, dashed] (0,-11.108) -- (1,-0.186) -- (2,22.070) -- (3,65.782);
\foreach \x/\y in {0/-11.108,1/0.906,2/27.836,3/86.016} {\fill[Rama4] (\x,\y) circle (3pt);}
\foreach \x/\y in {1/-0.186,2/22.070,3/65.782} {\fill[AzulUNA] (\x,\y) circle (3pt);}
\node[right, font=\normalsize\bfseries, text=Rama4] at (3.02,86.02) {86.02};
\node[right, font=\normalsize\bfseries, text=AzulUNA] at (3.02,65.78) {65.78};
\node[left, font=\normalsize\bfseries, text=Rama4] at (-0.02,-11.1) {$-$11.11};
\fill[RaciR] (0.92,0) circle (4pt); \fill[RaciR] (1.01,0) circle (4pt);
\draw[RaciR, line width=1.2pt, ->, >=Stealth] (0.62,22) node[above, font=\normalsize\bfseries, text=RaciR, align=center]{payback simple 0.92 años\\descontado 1.01 años} -- (0.92,1.5);
\end{tikzpicture}
\end{adjustbox}
{\normalsize \textcolor{Rama4}{\rule{18pt}{3pt}}~Flujo económico acumulado \quad \textcolor{AzulUNA}{\rule{18pt}{3pt}}~Acumulado descontado al 10\,\%}
\end{minipage}\hfill
\begin{minipage}[c]{9.3cm}
{\normalsize
\rowcolors{2}{white}{GrisFila}
\noindent\begin{tabularx}{\linewidth}{@{}L Q{2.9cm}@{}}
\toprule
\rowcolor{AzulClaro}
\textbf{Indicador} & \textbf{Resultado} \\
\midrule
Inversión (año 0) & S/ 11,108.07 \\
Flujo económico del año 1 & S/ 12,013.78 \\
\textbf{Payback simple} & \textbf{0.92 años} \\
Payback descontado & 1.01 años \\
VAN económico (10\,\%) & S/ +65,781.59 \\
TIR económica & 169.9\,\% \\
Relación beneficio/costo & 6.92 \\
ROI a 3 años & 774.4\,\% \\
VAN financiero & S/ $-$15,355.40 \\
\bottomrule
\end{tabularx}
}

\vspace{0.3cm}
{\normalsize\noindent\textbf{Cómo se calcula:} payback = inversión / flujo del año 1 = 11,108.07 / 12,013.78 = \textbf{0.92 años}, unos 11 meses. El flujo económico suma al flujo de la app el beneficio para los arrendatarios: S/ 41.64 por persona al año (350, 700 y 1,400 personas). Sin ese beneficio el VAN es negativo: es un proyecto de valor social.}
\end{minipage}
\end{hojah}

\begin{hojah}{Riesgos: Matriz de Probabilidad e Impacto (Versión 1.1)}
\begin{minipage}[c]{10.4cm}
\begin{adjustbox}{width=\linewidth}
\begin{tikzpicture}[font=\large]
\foreach \p in {1,2,3} {\foreach \i in {1,2,3} {
  \pgfmathtruncatemacro{\sv}{\p*\i}
  \ifnum\sv>5 \def\cc{SevAlta}\else\ifnum\sv>2 \def\cc{SevMedia}\else\def\cc{SevBaja}\fi\fi
  \fill[\cc!30] ({(\i-1)*3.2},{(\p-1)*2.6}) rectangle ({\i*3.2},{\p*2.6});
  \draw[white, line width=2pt] ({(\i-1)*3.2},{(\p-1)*2.6}) rectangle ({\i*3.2},{\p*2.6});
  \node[font=\small, text=GrisOscuro, anchor=north east] at ({\i*3.2-0.1},{\p*2.6-0.1}) {\sv};}}
\node[font=\Large\bfseries] at (8.0,6.5) {R02};
\node[font=\Large\bfseries, align=center] at (8.0,3.9) {R01\\R08};
\node[font=\Large\bfseries, align=center] at (4.8,3.9) {R05 R06\\R07};
\node[font=\Large\bfseries] at (8.0,1.3) {R03};
\node[font=\Large\bfseries] at (1.6,3.9) {R04};
\foreach \i/\t in {1/Bajo,2/Medio,3/Alto} {\node[font=\normalsize] at ({(\i-0.5)*3.2},-0.4) {\t};}
\foreach \p/\t in {1/Baja,2/Media,3/Alta} {\node[font=\normalsize, anchor=east] at (-0.1,{(\p-0.5)*2.6}) {\t};}
\node[font=\normalsize\bfseries] at (4.8,-1.0) {Impacto};
\node[font=\normalsize\bfseries, rotate=90] at (-1.7,3.9) {Probabilidad};
\end{tikzpicture}
\end{adjustbox}
\end{minipage}\hfill
\begin{minipage}[c]{12.4cm}
{\normalsize
\rowcolors{2}{white}{GrisFila}
\noindent\begin{tabularx}{\linewidth}{@{}C{1.0cm} L C{1.0cm} P{2.1cm}@{}}
\toprule
\rowcolor{AzulClaro}
\textbf{ID} & \textbf{Riesgo y respuesta} & \textbf{Sev.} & \textbf{Estado} \\
\midrule
R02 & Retraso de la nube: Ruth apoya la seguridad por fila & 9 & \etiqueta{SevAlta}{Activado} \\
R08 & Piloto fuera de la semana 13: cambio SC-02 & 6 & \etiqueta{SevAlta}{Nuevo} \\
R01 & Pocos propietarios: visitas de los sábados & 6 & \etiqueta{SevMedia}{Vigilancia} \\
R07 & Baja adopción: ya hay 14 inscritos & 4 & \etiqueta{SevMedia}{Vigilancia} \\
R05 & Ausencia de un integrante: apoyo por tarjeta & 4 & \etiqueta{Rama4}{Sin cambio} \\
R06 & Límites del plan gratuito: uso menor al 10\,\% & 4 & \etiqueta{Rama4}{Sin cambio} \\
R03 & Falla de un equipo: todo en GitHub & 3 & \etiqueta{Rama4}{Sin cambio} \\
R04 & Alza de precios: se acepta & 2 & \etiqueta{Rama4}{Sin cambio} \\
\bottomrule
\end{tabularx}
}

\vspace{0.2cm}
{\normalsize\noindent Severidad $=$ probabilidad $\times$ impacto (1 a 9). Contingencia $= \sum$ impacto en soles $\times$ probabilidad $=$ S/ 555.40, aún sin uso.}
\end{minipage}
\end{hojah}

\begin{hojah}{Diagrama de Flujo: Control de Cambios}
\begin{grande}
\begin{tikzpicture}[fl]
\node[term] (ini) {Necesidad de cambio};
\node[proc=3.4cm, right=of ini] (iss) {Registrar Issue con etiqueta \texttt{cambio}};
\node[proc=3.6cm, right=of iss] (ana) {Análisis de impacto en 3 días: alcance, tiempo, costo, calidad y riesgos};
\node[dec=2.6cm, right=of ana] (may) {Es un cambio mayor?};
\node[proc=3.4cm, above=1.2cm of may] (pm) {Decide el PM\\{\small hasta 2 días y S/ 50, sin tocar entregables}};
\node[proc=3.4cm, below=1.4cm of may] (vot) {Vota el sponsor en el control del martes};
\node[dec=2.4cm, right=1.4cm of vot] (apr) {Unanimidad?};
\node[proc=3.4cm, right=1.2cm of apr] (rech) {Rechazado: queda en el registro con su motivo};
\node[proc=3.4cm, above=1.4cm of rech] (act) {Actualizar líneas base, tablero y documentos (versión nueva)};
\node[term, above=1.0cm of act] (fin) {Comunicar en el acta};
\draw[fle] (ini) -- (iss);
\draw[fle] (iss) -- (ana);
\draw[fle] (ana) -- (may);
\draw[fle] (may) -- node[etq, right]{no} (pm);
\draw[fle] (may) -- node[etq, right]{sí} (vot);
\draw[fle] (vot) -- (apr);
\draw[fle] (apr) -- node[etq, above]{no} (rech);
\draw[fle] (apr) |- node[etq, pos=0.25, left]{sí} (act);
\draw[fle] (pm) -| (act);
\draw[fle] (act) -- (fin);
\draw[fle] (rech.east) -- ++(0.6,0) |- (fin.east);
\end{tikzpicture}
\end{grande}
\vspace{0.3cm}
{\normalsize
\rowcolors{2}{white}{GrisFila}
\noindent\begin{tabularx}{\linewidth}{@{}C{1.4cm} C{1.7cm} L P{3.6cm} C{2.4cm} C{2.4cm}@{}}
\toprule
\rowcolor{AzulClaro}
\textbf{ID} & \textbf{Fecha} & \textbf{Cambio o acción} & \textbf{Impacto principal} & \textbf{Responsable} & \textbf{Decisión} \\
\midrule
AC-01 & 29/09 & Ejecución en paralelo: el ranking avanza junto con la documentación & Cronograma (SPI 0.56) & Clyde & \etiqueta{Rama4}{Aplicada} \\
SC-01 & 02/10 & Pasar de 8 a 10 horas semanales por integrante & Costo y equipo & Clyde & \etiqueta{SevAlta}{Rechazada} \\
SC-02 & 10/11 & Mover el piloto a la semana 14 (23 al 27/11) & Cronograma (R08) & Ruth & \etiqueta{SevMedia}{En análisis} \\
SC-03 & 13/11 & Bajar la meta de 15 a 12 cuartos para el piloto & Alcance (R01) & Clyde & \etiqueta{SevMedia}{Por presentar} \\
\bottomrule
\end{tabularx}
}
\end{hojah}

\section{Calidad del Producto}
{\small
\rowcolors{2}{white}{GrisFila}
\noindent\begin{tabularx}{\textwidth}{@{}P{3.6cm} L C{2.4cm} C{2.2cm}@{}}
\toprule
\rowcolor{AzulClaro}
\textbf{ISO/IEC 25010:2023} & \textbf{Métrica} & \textbf{Meta} & \textbf{Hoy} \\
\midrule
Adecuación funcional & Requisitos aceptados y desglose correcto & 100\,\% & 6 de 9 \\
Eficiencia de desempeño & Tiempo del ranking con 500 cuartos & $\le$ 250 ms & Cumple \\
Capacidad de interacción & Tiempo para publicar sin ayuda; satisfacción & $\le$ 10 min; $\ge$ 80\,\% & Piloto \\
Fiabilidad & Errores críticos abiertos al entregar & 0 & Piloto \\
Seguridad & Accesos indebidos entre roles & 0 & Con RLS \\
Mantenibilidad & Pruebas automáticas del motor & 7 de 7 & 12 de 12 \\
Flexibilidad & Funciona en Android 10 y en uno reciente & 2 de 2 & Piloto \\
\bottomrule
\end{tabularx}
}

\vspace{0.1cm}
{\small\noindent Hoy pasan 12 pruebas automáticas del motor CASA SCORE y del cálculo de distancia. Cada Pull Request lo revisa un integrante distinto al autor y la web se instala desde cero con \texttt{git clone}, \texttt{npm install} y \texttt{npm run dev}.}

\Needspace{40\baselineskip}
\section{Lista de Verificación de Entregables}
\newcommand{\ok}{\textcolor{Rama4}{\large$\checkmark$}}
\newcommand{\pend}{\textcolor{GrisOscuro}{\large$\square$}}
{\footnotesize
\setlength{\tabcolsep}{3pt}
\rowcolors{2}{white}{GrisFila}
\begin{longtable}{@{}C{0.7cm} C{0.8cm} >{\raggedright\arraybackslash}p{6.2cm} P{2.6cm} C{1.5cm} P{3.4cm}@{}}
\toprule
\rowcolor{AzulClaro}
& \textbf{N.\textsuperscript{o}} & \textbf{Entregable} & \textbf{Responsable} & \textbf{Fecha} & \textbf{Evidencia en GitHub} \\
\midrule
\endfirsthead
\toprule
\rowcolor{AzulClaro}
& \textbf{N.\textsuperscript{o}} & \textbf{Entregable} & \textbf{Responsable} & \textbf{Fecha} & \textbf{Evidencia en GitHub} \\
\midrule
\endhead
\ok & C1 & Políticas y gobernanza del equipo & Viviana & 04/09 & PR \#10 \\
\ok & C2 & Acta de conformación del equipo & Viviana & 01/09 & PR \#9 \\
\ok & C6 & Acta de compromiso del equipo & Clyde & 09/09 & commit 3b78d13 \\
\ok & 01 & Ficha técnica & Clyde & 02/09 & commit c20617f \\
\ok & 02 & Especificación de requerimientos & Eliseo & 04/09 & PR \#14 \\
\ok & 03 & Viabilidad y caso de negocio & Ruth & 10/09 & PR \#1 \\
\ok & 04 & Acta de Constitución & Clyde & 11/09 & commit c20617f \\
\ok & 05 & Registro de interesados & Clyde & 18/09 & commit c20617f \\
\ok & 06 & EDT y diccionario & Ruth & 18/09 & PR \#2 \\
\ok & 07 & Plan del equipo & Ruth & 18/09 & PR \#3 \\
\ok & 08 & Cronograma & Ruth & 18/09 & PR \#4 \\
\ok & 09 & Presupuesto & Ruth & 18/09 & PR \#5 \\
\ok & 10 & Plan y registro de riesgos (versión 1.1) & Viviana & 22/09 y 06/11 & PR \#11 y commit fdbba86 \\
\ok & 11 & Plan de calidad & Eliseo & 22/09 & PR \#15 \\
\ok & 12 & Plan de comunicaciones & Viviana & 22/09 & PR \#12 \\
\ok & 13 & Plan de contratos & Eliseo & 22/09 & PR \#16 \\
\ok & 14 & Arquitectura técnica & Eliseo & 22/09 & PR \#17 \\
\ok & 15 & Actas de los sprints 1 a 5 & Viviana & 16/10 & PR \#13 \\
\ok & 16 & Control del proyecto & Ruth & 02/10 & PR \#7 \\
\ok & 17 & Producto e informe de validación & Eliseo & 09/10 & PR \#18 y \#19 \\
\ok & 18 & Sustentación de la Unidad I & Clyde & 16/10 & commit 2ea9619 \\
\ok & 19 & Guía de marcos de referencia & Clyde & 20/10 & commit 2ea9619 \\
\ok & 20 & Plan de Dirección y su resumen & Clyde & 26/10 & commit 2ea9619 \\
\ok & 21 & Informe de desempeño & Ruth & 06/11 & PR \#8 \\
\ok & -- & App web con ranking, mapa y publicación & Eliseo & 06/11 & PR \#20 \\
\ok & 00 & Expediente del proyecto (resumen gráfico) & Clyde & 09/11 & carpeta 00\_expediente \\
\ok & 29 & Bosquejo del cierre del proyecto & Clyde & 13/11 & carpeta 29\_bosquejo\_cierre \\
\pend & 22 & Registro de adaptación & Clyde & 13/11 & -- \\
\pend & 23 & Plan de cierre y lecciones aprendidas & Viviana & 13/11 & -- \\
\pend & 24 & Informe de auditoría y calidad & Eliseo & 27/11 & -- \\
\pend & 25 & Informe de integración de procesos & Clyde & 27/11 & -- \\
\pend & 26 & Manuales y versión 1.0 del MVP & Eliseo & 04/12 & -- \\
\pend & 27 & Informe final y exposición & Clyde & 14/12 & -- \\
\pend & 28 & Acta de cierre & Clyde & 18/12 & -- \\
\bottomrule
\end{longtable}
}

\Needspace{16\baselineskip}
\section{Lista de Verificación de Gates}
{\small
\rowcolors{2}{white}{GrisFila}
\noindent\begin{tabularx}{\textwidth}{@{}C{0.7cm} C{0.9cm} L C{1.8cm}@{}}
\toprule
\rowcolor{AzulClaro}
& \textbf{Gate} & \textbf{Condición para pasar} & \textbf{Fecha} \\
\midrule
\ok & G1 & Acta de Constitución aprobada con objetivos medibles y líneas base iniciales. & 11/09 \\
\ok & G2 & MVP de la Unidad I con 5 cuartos en el orden esperado y 7 pruebas aprobadas. & 09/10 \\
\pend & G3 & Registro en la nube funcionando y piloto realizado con 10 arrendatarios y 5 propietarios. & 27/11 \\
\pend & G4 & RF-01 a RF-09 aceptados, 0 errores críticos y 14 casos de prueba aprobados. & 04/12 \\
\pend & G5 & Informe final expuesto, lecciones aprendidas y acta de cierre firmada. & 18/12 \\
\bottomrule
\end{tabularx}
}

\vspace{0.15cm}
{\small\noindent\textbf{Definición de terminado:} una tarea pasa a Hecho cuando cumple su criterio de aceptación, otro integrante la revisó, su evidencia (commit o Pull Request) está enlazada en la tarjeta y sus horas reales están registradas.}

\section{Cierre de la Fase 1}
El proyecto del semestre es la fase 1 de un producto pensado para tres años o más. En diciembre se cierra la fase con el MVP funcionando en la nube, probado con cuartos reales de compañeros, y el producto pasa a operación. El detalle está en el bosquejo del cierre del proyecto.
\begin{hojah}{Proceso de Cierre de la Fase (PMBOK)}
\begin{grande}[12.6cm]
\begin{tikzpicture}[>=Stealth]
\newcommand{\paso}[6]{
\node[draw=#4, fill=#4!10, line width=1.3pt, rounded corners=6pt, minimum width=5.6cm, minimum height=4.4cm] (#1) at (#2) {};
\node[circle, fill=#4, text=white, font=\Large\bfseries, minimum size=0.9cm] at ($(#2)+(-2.25,1.7)$) {#3};
\node[font=\large\bfseries, text=#4, anchor=west, text width=4.3cm] at ($(#2)+(-1.7,1.7)$) {#5};
\node[font=\large, text width=5.1cm, align=left] at ($(#2)+(0,-0.4)$) {#6};}
\paso{a}{0,0}{1}{AzulUNA}{Verificar\\entregables}{Cada entregable contra su criterio de aceptación. \textbf{Hoy:} 21 de 28 documentos y MVP en la nube.}
\paso{b}{6.0,0}{2}{AzulUNA}{Aceptación formal}{Votación del sponsor en el gate G5 y acta de cierre firmada. \textbf{Fecha:} 18/12.}
\paso{c}{12.0,0}{3}{AzulUNA}{Cerrar contratos}{Términos de los planes gratuitos revisados y cuenta de Google Play definida. \textbf{Fecha:} 11/12.}
\paso{d}{18.0,0}{4}{AzulUNA}{Liberar recursos}{Horas del semestre, celular de pruebas y accesos de cada integrante. \textbf{Fecha:} 18/12.}
\paso{e}{18.0,-5.2}{5}{Rama4}{Archivar}{Líneas base, cambios, actas e informes en el repositorio con su versión final. \textbf{Fecha:} 16/12.}
\paso{f}{12.0,-5.2}{6}{Rama4}{Lecciones\\aprendidas}{Lo que funcionó, lo que no y qué haremos distinto, por rol. \textbf{Primer bosquejo:} hoy.}
\paso{g}{6.0,-5.2}{7}{Rama4}{Comunicar\\el cierre}{Informe final y exposición ante el asesor y los clientes del piloto. \textbf{Fecha:} 14/12.}
\paso{h}{0,-5.2}{8}{Rama4}{Medir beneficios}{Satisfacción del piloto ahora y beneficio por arrendatario desde 2027.}
\foreach \x/\y in {a/b,b/c,c/d,e/f,f/g,g/h} {\draw[->, line width=2pt, AzulUNA] (\x) -- (\y);}
\draw[->, line width=2pt, AzulUNA] (d.south) -- (e.north);
\end{tikzpicture}
\end{grande}
\vspace{0.2cm}
{\large\noindent Los pasos 1 a 4 cierran lo administrativo; los pasos 5 a 8 conservan el conocimiento y el valor para las fases siguientes.}
\end{hojah}
\begin{hojah}{Ciclo de Vida del Software: Despliegue y Pruebas Continuas en Producción}
\begin{grande}[11.8cm]
\begin{tikzpicture}[>=Stealth]
\newcommand{\et}[5]{\node[draw=#3, fill=#3!#5, line width=1.3pt, rounded corners=6pt, minimum width=3.6cm, minimum height=2.6cm, align=center, font=\large\bfseries, text=#4] (#1) at (#2) {};}
\et{r}{0,0}{GrisOscuro}{black}{10}
\et{d}{4.2,0}{GrisOscuro}{black}{10}
\et{c}{8.4,0}{GrisOscuro}{black}{10}
\et{p}{12.6,0}{GrisOscuro}{black}{10}
\et{dep}{16.8,0}{RaciR}{RaciR}{12}
\et{pc}{21.0,0}{RaciR}{RaciR}{12}
\foreach \n/\t/\s in {r/Requisitos/RF-01 a RF-09, d/Diseño/{arquitectura\\y pantallas}, c/Construcción/sprints semanales, p/Pruebas/{12 automáticas\\y 14 casos}, dep/Despliegue/MVP en la nube, pc/{Pruebas\\continuas}/en producción} {
  \node[font=\large\bfseries, align=center] at ($(\n)+(0,0.5)$) {\t};
  \node[font=\normalsize, align=center] at ($(\n)+(0,-0.6)$) {\s};}
\foreach \a/\b in {r/d,d/c,c/p,p/dep,dep/pc} {\draw[->, line width=2pt, AzulUNA] (\a) -- (\b);}
\draw[->, line width=2pt, Rama4, dashed] (pc.south) -- ++(0,-0.9) -| node[pos=0.25, below, font=\large\bfseries, text=Rama4]{retroalimentación de los usuarios al backlog} (c.south);
\node[font=\Large\bfseries, text=RaciR] at (18.9,1.75) {Etapas de esta fase};
\draw[RaciR, line width=1.2pt, dashed, rounded corners=8pt] (14.8,-1.4) rectangle (23.0,2.15);
% detalle despliegue
\node[draw=RaciR, fill=white, line width=1.1pt, rounded corners=5pt, text width=10.6cm, align=left, font=\large, anchor=north west] at (0,-3.3) {\textbf{\textcolor{RaciR}{Despliegue}}\\[2pt]
1. Pull Request revisado por otro integrante y unido a \texttt{main}.\\
2. Pasan las 12 pruebas automáticas y la compilación sin errores.\\
3. Se publica la versión en la nube y se etiqueta (v0.x hasta la v1.0).\\
4. Prueba de humo: ingresar, ver el ranking, abrir un detalle y publicar.\\
5. Si algo falla, se vuelve a la versión anterior en el mismo día.};
\node[draw=RaciR, fill=white, line width=1.1pt, rounded corners=5pt, text width=11.4cm, align=left, font=\large, anchor=north west] at (11.4,-3.3) {\textbf{\textcolor{RaciR}{Pruebas continuas en producción}}\\[2pt]
1. Cuartos reales de compañeros publicados y revisados cada semana.\\
2. Ranking desde distintos puntos de Puno: orden y desglose correctos.\\
3. Tiempo de respuesta del ranking menor a 250 ms.\\
4. Uso de los planes gratuitos menor al 70\,\% de sus límites.\\
5. Cada defecto se registra como tarjeta \textit{Defecto} en el tablero.};
\end{tikzpicture}
\end{grande}
\end{hojah}
\section{Revisión de Riesgos Cada Dos Semanas}
Desde la semana 12 el registro de riesgos se revisa a fondo cada dos semanas, además del control semanal de los disparadores. En cada revisión se recalifican todos los riesgos, se cierran los que pasaron su ventana y se agregan los del cierre.

\vspace{0.1cm}
\begin{center}
\begin{tikzpicture}[font=\small]
\draw[line width=1.5pt, AzulUNA] (0,0) -- (15.6,0);
\foreach \x/\d/\t in {0.6/13-11/{Revisión 1\\riesgos del cierre}, 4.6/27-11/{Revisión 2\\tras el piloto}, 8.6/11-12/{Revisión 3\\antes del informe}, 12.6/18-12/{Revisión final\\riesgos al año 1}} {
  \node[circle, fill=RaciR, minimum size=0.35cm] at (\x,0) {};
  \node[above=4pt, font=\small\bfseries, text=RaciR] at (\x,0.15) {\d};
  \node[below=6pt, align=center] at (\x,-0.1) {\t};}
\end{tikzpicture}
\end{center}

{\small
\rowcolors{2}{white}{GrisFila}
\noindent\begin{tabularx}{\textwidth}{@{}C{0.9cm} L C{0.6cm} C{0.6cm} C{0.9cm} P{3.6cm}@{}}
\toprule
\rowcolor{AzulClaro}
\textbf{ID} & \textbf{Riesgo del cierre} & \textbf{P} & \textbf{I} & \textbf{Sev.} & \textbf{Respuesta} \\
\midrule
R02 & El registro en la nube no queda listo para el piloto. & 3 & 3 & \etiqueta{SevAlta}{\,9\,} & Apoyo de Ruth y piloto en la semana 14. \\
R09 & Los clientes del piloto no aceptan el MVP como útil. & 2 & 3 & \etiqueta{SevAlta}{\,6\,} & Guion de prueba y encuesta con criterios claros. \\
R10 & El MVP en la nube se cae o supera los límites gratuitos. & 2 & 2 & \etiqueta{SevMedia}{\,4\,} & Monitoreo semanal y vuelta a la versión anterior. \\
R11 & Las lecciones aprendidas quedan incompletas por los exámenes. & 2 & 2 & \etiqueta{SevMedia}{\,4\,} & Plantilla por rol entregada el 13/11. \\
R12 & No se define quién opera el producto en el año 1. & 2 & 3 & \etiqueta{SevAlta}{\,6\,} & Acuerdo de continuidad en el acta de cierre. \\
\bottomrule
\end{tabularx}
}


\Needspace{12\baselineskip}
\section{Lecciones Principales Hasta Hoy}
{\small
\rowcolors{2}{white}{GrisFila}
\noindent\begin{tabularx}{\textwidth}{@{}P{3.0cm} L L@{}}
\toprule
\rowcolor{AzulClaro}
& \textbf{Qué ocurrió} & \textbf{Qué haremos distinto} \\
\midrule
\textcolor{Rama4}{\textbf{Funcionó}} & Motor separado y probado; MVP de la Unidad I adelantado. & Separar la lógica desde el primer sprint. \\
\textcolor{Rama4}{\textbf{Funcionó}} & Tablero con evidencia obligatoria y demostración cada viernes. & Mantenerlo en la operación del producto. \\
\textcolor{RaciR}{\textbf{Mejorar}} & Planificación firmada con dos semanas de retraso (SPI 0.56). & Planificación mínima firmada antes de construir. \\
\textcolor{RaciR}{\textbf{Mejorar}} & Sin prueba de concepto de la nube; R02 activado. & Prueba técnica corta en la primera iteración. \\
\bottomrule
\end{tabularx}
}

\begin{hojah}{Perspectiva del Producto: Hoja de Ruta a Tres Años}
\begin{grande}[12.4cm]
\begin{tikzpicture}[>=Stealth]
\draw[line width=3pt, AzulUNA, ->] (0,0) -- (24.2,0);
\foreach \x/\t/\c/\s/\d in {0.3/2026-II: fase 1/RaciR/{Proyecto del semestre}/{MVP en la nube, piloto con compañeros y clientes, expediente y cierre}, 6.3/2027: año 1/Rama4/{Operación en Puno}/{500 usuarios al mes, app en Google Play, suscripción de S/ 5 por cuarto y publicidad local y en internet}, 12.3/2028: año 2/Rama2/{Crecimiento}/{1,000 usuarios al mes, Juliaca, calificaciones y reportes de publicaciones}, 18.3/2029: año 3/Rama5/{Expansión}/{2,000 usuarios al mes y llegada a Cusco y Arequipa}} {
  \node[circle, fill=\c, minimum size=0.6cm] at (\x+2.6,0) {};
  \node[fill=\c, text=white, font=\large\bfseries, minimum width=5.4cm, minimum height=0.9cm, anchor=south west] at (\x,0.6) {\t};
  \node[draw=\c, fill=\c!8, line width=1.2pt, text width=5.0cm, align=left, font=\normalsize, anchor=north west, minimum height=4.0cm] at (\x,-0.6) {\textbf{\s}\\[3pt]\d};}
\node[font=\large\bfseries, text=RaciR, align=center] at (2.9,3.0) {Se cierra en diciembre};
\draw[RaciR, line width=1.4pt, dashed] (5.95,-5.0) -- (5.95,3.3);
\node[font=\normalsize, text width=23cm, align=left, anchor=north west] at (0,-5.3) {Beneficio esperado para los arrendatarios: S/ 14,574 (año 1), S/ 29,148 (año 2) y S/ 58,296 (año 3). VAN económico de S/ 65,781.59 y payback de 0.92 años desde el inicio de la operación.};
\end{tikzpicture}
\end{grande}
\end{hojah}


\end{document}
